\documentclass[11pt,a4paper]{article}
\usepackage[margin=25mm]{geometry}
\usepackage{amsmath,amssymb,amsthm,mathtools}
\usepackage{booktabs,longtable,array,enumitem}
\usepackage[expansion=false]{microtype}
\usepackage{xcolor}
\usepackage[colorlinks=true,linkcolor=blue!45!black,citecolor=blue!45!black,urlcolor=blue!45!black]{hyperref}
\usepackage{fancyhdr}
\pdfmapfile{+cm.map}
\pdfmapfile{+cmextra.map}
\pdfmapfile{+symbols.map}
\pdfmapfile{+latxfont.map}
\setlength{\emergencystretch}{3em}
\setlength{\parskip}{4pt}
\setlength{\parindent}{0pt}
\setlist{itemsep=3pt,topsep=5pt,leftmargin=*}
\pagestyle{fancy}\fancyhf{}\fancyhead[L]{\small GCT\_NSRH: checked reductions and open gates}\fancyfoot[C]{\thepage}
\setlength{\headheight}{14pt}
\newcommand{\Q}{\mathbb Q}\newcommand{\Z}{\mathbb Z}\newcommand{\C}{\mathbb C}\newcommand{\N}{\mathbb N}\newcommand{\F}{\mathbb F}
\newcommand{\A}{\mathbb A}\newcommand{\PP}{\mathbb P}\newcommand{\GL}{\mathrm{GL}}
\newcommand{\Hom}{\operatorname{Hom}}\newcommand{\Spec}{\operatorname{Spec}}\newcommand{\Tr}{\operatorname{Tr}}\newcommand{\rank}{\operatorname{rank}}\newcommand{\Sym}{\operatorname{Sym}}\newcommand{\Res}{\operatorname{Res}}\newcommand{\IC}{\operatorname{IC}}\newcommand{\Gr}{\operatorname{Gr}}\newcommand{\Rep}{\operatorname{Rep}}\newcommand{\Frob}{\operatorname{Frob}}\newcommand{\End}{\operatorname{End}}
\newcommand{\id}{\operatorname{id}}\newcommand{\et}{\mathrm{\acute et}}\newcommand{\boxt}{\boxtimes}
\newcommand{\mathscr}{\mathcal}
\renewcommand{\S}{\ensuremath{\mathsection}}
\renewcommand{\textbullet}{\ensuremath{\bullet}}
\newcommand{\fid}[1]{\textsf{#1}}\newcommand{\gate}[1]{\textsf{#1}}
\newtheoremstyle{named}{8pt}{5pt}{\itshape}{} {\bfseries}{.}{.5em}{}
\theoremstyle{named}\newtheorem{finding}{Finding}[section]
\theoremstyle{definition}\newtheorem{definition}[finding]{Definition}
\theoremstyle{remark}\newtheorem{remark}[finding]{Remark}
\title{Purity from path sums and the nonstandard GCT positivity problem\\[4pt]\large Checked reductions, explicit barriers, and the missing realization}
\author{Research run \texttt{GCT\_NSRH\_20261121}}
\date{21 November 2026 (Europe/Berlin)}
\begin{document}\maketitle
\begin{abstract}
We investigate whether the path-quadric sign mechanism described in Alp\"oge's supplied \emph{Remark} transfers to Mulmuley's nonstandard quantum groups. We prove an exact multidimensional path-sum estimate and a conditional purity criterion. For trace functions on $\A^d/\F_q$ bounded in pointwise weight $a$, iterated Parseval gives an error of size $q^{n(k(a+d)/2+d/2+1)}$. If the relevant spectra have opposite cohomological parity for unbounded path lengths $k$, positive spectral multiplicities prevent cancellation. Taking first extension degree $n$ and then $k$ to infinity gives the upper weight $a+d$; a specified Tate duality gives equality. An explicit sublinear boundary-error version is included.

The supplied GCT papers do not formulate a finite-field nonstandard Riemann hypothesis with designated Frobenius objects and weights. Their positivity targets also include signed factors, rational normalizations, and, in general plethysm, corrections to positivity. We isolate the additional realization, effective-category, basis, Lefschetz, specialization, and complexity hypotheses. A rank-one calculation in the actual nonstandard Hecke inclusion proves that its ambient Kazhdan--Lusztig positive cone meets the nonstandard algebra only in scalar multiples of the unit. Thus a straightforward inheritance of the standard effective category fails even in this small case. Pure two-point geometry supplies a separate counterexample to determinant positivity implying positive structure constants in a chosen effective basis.

Theta/interlacing arguments yield an exact conditional path-count theorem, but no functor identifying those paths with the required GCT multiplicity spaces is supplied. Existing rectangular positivity and no-occurrence theorems are stated with their parameters. All checked findings are translated into schemes, representation $K_0$, $\lambda$-operations and Weil realizations, with verified archival anchors.
\end{abstract}
\begin{center}\fbox{\parbox{.91\linewidth}{\textbf{Research status: INCOMPLETE; work in progress.} The conditional theorems and counterexamples below are proved and checked. No nonstandard RH for the designated GCT objects, general canonical-basis positivity, general positive plethysm rule, or complexity separation is proved. The final part records the exact open gates.}}
\end{center}
\newpage\setcounter{tocdepth}{1}{\small\tableofcontents}\newpage

\part{The exact question and its sources}
\section{Scope, premises, and conventions}\label{sec:scope}
The main question is a transfer question. We accept the RH and Dirichlet-GRH results in the supplied Alp\"oge material as stipulated premises, as requested. We do not re-referee those results. Their acceptance does not supply a theorem relating an automorphic or finite-field trace system to a nonstandard quantum canonical basis.

The relevant finite-field description is in \cite[printed p.~3, footnote~3, attached to the paragraph before the timeline; also step~3]{Alpoge}. Timeline step~5 itself concerns a von Mangoldt functional and continues onto p.~4. We use ``the step-5 mechanism'' in the prompt's sense, but retain this more precise locator. The note names a path-quadric identity, middle-cohomological signs, iterated Parseval, and the absence of cancellation of inverse roots. It does not state a standalone theorem with all the geometric hypotheses used below.

The supplied \texttt{ALPOGE\_Main.lean} consists of five exported wrappers, including the RH wrapper \texttt{riemannHypothesis}. Its body is \texttt{rhOfAll}; the imported module is \texttt{Supplements.Kernels.\allowbreak MainTheoremsAll}. The supplied \texttt{ALPOGE\_Statements.lean} defines the zeta/Dirichlet zero and nonvanishing propositions. Neither file contains a nonstandard quantum-group, path-concentration, or theta/GCT transfer theorem. The imported development is not in this ZIP. The supplied scan records zero occurrences for the listed prohibited constructs; this run did not repeat that scan or execute the missing development. These are access and locator facts, not objections to the stipulated theorem. No Masullo result is used: none is identified or supplied in this prompt.

We write $v$ for the \emph{formal quantum parameter}, $q$ for a fixed prime power, and $Q=q^n$ for an extension-field cardinality. In the GCT sources the quantum parameter is usually called $q$. No substitution $v=Q$ or $v=Q^{1/2}$ constructs a Weil structure. All weights in the purity conclusion are integers. The Frobenius convention is geometric Frobenius, so $\Q_\ell(-1)$ has eigenvalue $q$ over $\F_q$.

``Finding'' means a deduction established for this investigation, including elementary lemmas in the precise form needed. It is not a claim of priority over the mathematical literature. Known research results are cited, not presented as new findings. A source statement checked in context is not thereby an independently checked proof of the entire cited theorem.

\section{What IV, VI, VII and VIII actually specify}\label{sec:sources}
\subsection{Versions and the missing RH statement}
The attached IV is the revised Blasiak--Mulmuley--Sohoni paper, arXiv cs/0703110v4 (2013), not merely the 2007 version referred to by VII and VIII. It proves nonstandard Schur--Weyl duality and the two-row crystal result. VI is arXiv 0704.0229v4 (2009); VII and VIII are 0709.0749v2 and 0709.0751v2 (2008). Rendered title-page dates produced by an undated \verb|\maketitle| are not used as version dates.

\textbf{Source finding A01.} In these four supplied versions there is no formulated nonstandard RH specifying finite-field objects, Frobenius actions, cohomological degrees and weights. This is an exact-source conclusion, not a theorem about what other work might contain. The relevant locators are:

{\small\begin{longtable}{p{.13\linewidth}p{.35\linewidth}p{.44\linewidth}}
\toprule Work&Exact locator&Content and status\\\midrule\endhead
IV&\S1.1, Problem 1.1; \S\S1.2--1.5; \S19&Positive-rule and basis program; the revised text has no finite-field purity statement.\\
VI&Introduction, Figure 1.1; \S1.7, Observation 1.7.1 and Figure 1.2&The proposed RH/extensions connection is marked by question marks; standard canonical positivity motivates the program.\\
VI&\S4.2, Proposition 4.2.1; \S4.2.2, Question 4.2.6 (PH0??); \S4.3&Asks for compatible canonical bases and analogous positivity. No nonstandard weight theorem is stated.\\
VII&\S1 after Conjectures 1.3--1.4; \S4.1, equations (9)--(10), Theorem 4.4; \S4.2, Theorem 4.5&Nonstandard minors and Laplace coefficients need interpretations analogous to perverse-sheaf constructions.\\
VIII&\S1, pp.~2--3; \S2.3 after Conjecture 2.5&An extension of perverse sheaves and finite-field RH is proposed as a way to understand minors and globalization.\\
VIII&\S2.4.2 after Conjecture 2.7; \S2.4.3 after Conjecture 2.9&Topological interpretations, Hard Lefschetz, and a theory for noncommutative, nonflat objects are suggested; the general case includes correction terms.\\\bottomrule
\end{longtable}}
These statements are supported by complete reading of their contexts and matching numbered PDF statements, not merely a search for the phrase ``Riemann hypothesis.'' Their appropriate conclusion is that a precise realization/purity package must first be specified.

\subsection{The algebraic objects}
Let $H$ be a complex connected classical reductive group, $X=V_\mu(H)$, $G=\GL(X)$, and $\rho:H\to G$ the representation homomorphism. The generalized plethysm multiplicity is
\begin{equation}\label{eq:plethysm}
a_{\lambda,\mu}^{\pi}=\dim\Hom_H\bigl(V_\pi,\Res_\rho V_\lambda(G)\bigr).
\end{equation}
For Kronecker coefficients, $H=\GL(V)\times\GL(W)$, $X=V\otimes W$, and $\rho(g,h)=g\otimes h$.

VII, Problem~1.1 and Theorems~1.2, 2.1 and 3.1, constructs a nonstandard matrix-coordinate bialgebra
\[
 S=\mathcal O(M_v^H(X)),\qquad \mathcal A_r^H=(S_r)^*,
\]
with a quantized map from $H_v$ to the associated nonstandard quantum object. Its matrix generators satisfy the projector relation
\[
 P^{+,H}_{X,X}(\mathbf u\otimes\mathbf u)
 =(\mathbf u\otimes\mathbf u)P^{+,H}_{X,X}
\]
(VII, \S2, equations (4)--(6)). The projector is constructed from the braiding on $X_v\otimes X_v$. The determinant may vanish, so the strict term may be quantum semigroup; an antipode is not automatic. The dual nonstandard symmetric-group deformation is $\mathcal B_r^H(v)$ (VII, \S5).

These are noncommutative algebras over fields containing $\Q(v)$, sometimes an algebraic extension. Their bar involution is not arithmetic Frobenius. Complete reducibility or a compact real form is not a Weil structure. Nonflatness and noncommutativity have different roles: nonflat commutative schemes still have \(\ell\)-adic sheaves, whereas a noncommutative coordinate algebra is not the coordinate ring of a scheme in the EGA/SGA sense.

General plethysm duality and reciprocity are VII, Conjectures~1.3--1.4 and refined Conjecture~6.1. They would give
\begin{equation}\label{eq:comparison}
a_{\lambda,\mu}^{\pi}=\sum_\alpha m_\lambda^\alpha n_\pi^\alpha,
\end{equation}
where $n_\pi^\alpha$ is an $H_v$ constituent multiplicity and $m_\lambda^\alpha$ is a Specht multiplicity after specialization at $v=1$. In revised IV, Theorems~1.5/12.1 establish the Kronecker duality, and \S1.4, equation~(11), gives the corresponding Kronecker decomposition. It would be incorrect to leave that particular comparison conjectural merely because VII predates the revision.

\subsection{The actual positivity targets}\label{sec:targets}
The basis constructions themselves have separate prerequisites. VIII, Conjecture~2.2, concerns compatible local crystal bases; Question~2.3's ordinary balanced-triple condition fails in examples; Conjecture~2.5 proposes a unique minimum-degree globalization. Conjectures~2.6 and 3.3 concern cells, Conjectures~3.1--3.2 the dual globalization, and Conjecture~3.7 the \emph{opposite} basis's quasi-subcells and specialization. VIII's Proposition~2.4 contains an omitted proof and assumes the local-crystal conjecture; we do not record it as independently proved.

\paragraph{Kronecker, coordinate side.}
VIII, Conjecture~2.7, asks that expansion coefficients of $CB(\mathcal A_r^H)$ in the external canonical basis $CB(Z)$, $Z=\End(X_v^{\otimes r})$, be polynomials with nonnegative \emph{rational} coefficients (equation~(9)). Multiplication and comultiplication constants have the form
\begin{equation}\label{eq:kroneckerpositive}
 \pm(v-v^{-1})^a f(v),\qquad a\in\Z_{\ge0},
\end{equation}
where $f$ is bar-invariant, coefficientwise nonnegative rational and unimodal. The same is proposed for $CB(S)$.

\paragraph{Kronecker, dual side.}
VIII, Conjecture~3.4, applies to structure constants and external expansion coefficients of $E(r)$ and $E^{\mathrm{opp}}(r)$. Its form is
\begin{equation}\label{eq:dualpositive}
 \pm(v^{1/2}-v^{-1/2})^a f(v),
\end{equation}
with $f$ nonnegative rational, bar-invariant and unimodal in the half-power variable. Thus raw coefficientwise positivity is not the stated assertion even in the Kronecker case.

\paragraph{General plethysm.}
VIII, Conjectures~2.8 and 3.5, concern normalized coefficients of minimal polynomials of algebraic structure/expansion coefficients. The displayed normalizations are equations~(10) and (14):
\[
 (-1)^e(v-v^{-1})^{e'}c(v),\qquad
 (-1)^e(v^{1/2}-v^{-1/2})^{e'}g(v).
\]
The nonzero normalized polynomial is required to be positive at $v=1$ (called saturated there). Conjectures~2.9 and 3.6 ask for a dominant positive unimodal part and a small correction, for example
\[
 c=c^0+c^1,\qquad
 \frac{\|c^1\|_2}{\|c^0\|_2}
 \le\frac1{\operatorname{poly}(\langle\mu\rangle,\langle\pi\rangle,\langle r\rangle)},
 \qquad c^0(1)\gg |c^1(1)|.
\]
Conjecture~3.6 is explicitly informal. A theorem target must fix the quantitative meaning of the polynomial and the asymptotic comparison. VIII, \S5.4.1, Figure~30, displays negative entries in a normalized example; these are cited source computations, not a recomputation of the whole algebra in this run.

There is also a sign-convention issue in ``coefficients of the minimal polynomial'': the usual monic polynomial of 1 is $T-1$, while the source later refers to trace and norm without the conventional sign on trace. We record this ambiguity rather than use it to disprove the intended conjecture. ``Frobenius automorphisms'' in VIII, \S5.4.1, refers to algebraic conjugations over $\Q(v)$, not to finite-field Weil Frobenius.

\subsection{Positivity hypotheses in the flip}
VI, Hypothesis~1.6.4 (PH1), is substantially stronger than numerical nonnegativity. For each triple there should be a rational polytope $P_{\lambda,\mu}^{\pi}$ whose Ehrhart function is
\[
 \widetilde a_{\lambda,\mu}^{\pi}(t)=a_{t\lambda,\mu}^{t\pi},
\]
with $\mu$ fixed (VI, \S1.6, definitions preceding Theorem~1.6.1). It must admit a description $Ax\le b$ with $A$ independent of $\lambda,\pi$ and $b$ homogeneous linear in them; ambient dimension and entry bitlengths are polynomially bounded as specified there; membership and a separating hyperplane are computable in polynomial time. The number of inequalities need not be polynomial.

VI, Hypothesis~1.6.5 (SH), asks for a polynomial bound on the saturation index in $\dim G+\operatorname{ht}\lambda+\operatorname{ht}\pi$ and asymptotic generic strict saturation. Hypothesis~1.6.6 (PH2) bounds the positivity index; Hypothesis~1.6.7 (PH3) bounds the modular index of a positive rational Hilbert-series form. Strict saturation in every case is not the attached revised hypothesis.

The conditional implication to relaxed PHflip is VI, Theorem~1.6.8. A printed discrepancy must be retained: its part~(a) displays $a_{b\lambda,b\mu}^{b\pi}$ although the stretching definition and its stated proof through PH1 keep $\mu$ fixed. We use only the fixed-$\mu$ consequence directly justified by the displayed Ehrhart model and record the all-three-scaling formulation as requiring clarification. This editorial point is not needed for our main obstruction. More fundamentally, VIII, \S6, explicitly asks for additional effective-label properties. An abstract canonical basis or an abstract positive count does not prove PH1, SH, or the existence of a lower-bound obstruction.

\section{The Kazhdan--Lusztig template, with weights fixed}\label{sec:KL}
Let $G/\F_Q$ be split reductive, $B$ a Borel subgroup, $C_w=BwB/B$ and $X_w=\overline{C_w}$. Put
\[
 \IC_w^u=j_{!*}\Q_\ell[\ell(w)][-\ell(w)]
\]
with the usual extension to $X_w$; the notation records the unshifted convention, constant in degree zero on $C_w$. Then
\begin{equation}\label{eq:KL}
 P_{x,w}(T)=\sum_{j\ge0}\dim\mathcal H^{2j}(\IC_w^u)_x\,T^j.
\end{equation}
Odd stalk groups vanish, and each eigenvalue in degree $2j$ is $Q^j$. This is the Tate/parity statement, stronger than an absolute-value bound. Exact primary locators for this form are Haines \cite[Theorems 2.0.1, 5.0.4; Proposition 3.0.2; \S4]{Haines}. The original KL geometric identification is \cite[Theorem 4.3]{KL}.

In the modern proof, a Bott--Samelson resolution is proper, its fibers admit affine pavings, and the decomposition theorem isolates the intersection complex. The affine paving provides even/Tate cohomology. The compatibility with Weil structures must be checked: a geometric splitting need not itself be Frobenius-equivariant. Haines \S4 discusses precisely this issue. His introduction also notes that the original KL proof does not require the decomposition theorem; the latter is the requested modern template, not a claim about historical logical necessity.

Deligne's upper-weight theorem is \cite[Theorem 3.3.1; Corollaries 3.3.4--3.3.6]{WeilII}. BBD intermediate extension preserves purity (5.3.2), pure perverse sheaves are geometrically semisimple (5.3.8), and pure complexes split geometrically (5.4.5); relative Hard Lefschetz is 5.4.10 \cite{BBD}. None of these statements identifies an arbitrary pure trace with the dimension polynomial \eqref{eq:KL}. Arithmetic Frobenius semisimplicity is distinct: BBD, pp.~138--139, explicitly discusses pure unipotent Jordan-block sheaves.

{\small\begin{longtable}{p{.25\linewidth}p{.29\linewidth}p{.38\linewidth}}
\toprule Standard ingredient&Nonstandard datum supplied&Unfilled role\\\midrule\endhead
Flag schemes, Schubert strata&$H_v$, nonstandard FRT bialgebra; support-defined cells&An actual space/category with those strata and stalk functors.\\
Distinguished $\IC_w$ and its normalized class&Candidate minimum-degree bases&Objects whose classes are the \emph{specified} bases; existence and uniqueness.\\
Weil structure and closed-point Frobenius&Algebraic bar involution&Arithmetic models, traces, degrees and integer weights.\\
Proper resolutions/convolution&Algebra multiplication and comultiplication&Correspondences or functors with a decomposition theorem.\\
Tate parity and effective multiplicities&Signed rational forms in VIII&Positive multiplicity objects, signs, factors and denominators.\\
Verdier duality and Hard Lefschetz&Bar symmetry/unimodality conjectures&Compatible duality and Lefschetz isomorphisms.\\
Crystal/path labels&Conjectural cells; proved two-row tableaux&Effective labels, specialization and polynomial-time verification in general.\\\bottomrule
\end{longtable}}
For standard quantum groups, quiver and flag constructions supply related geometric realizations; it would be misleading to assert that every quantum canonical-basis constant is a stalk of one Schubert variety.

\part{Checked findings}
\section{The sign mechanism and its exact strength}\label{sec:main}
This section contains the main proved reduction. Its hypotheses explicitly retain the information that the phrase ``middle cohomology'' compresses.

\begin{finding}[\fid{F01}: path identity and Parseval estimate]\label{f:path}
Let $V=\F_Q^d$, $d\ge1$, with a nondegenerate bilinear pairing, and let $\psi$ be a nontrivial additive character. For $k\ge2$ and any function $t:V\to\C$, put
\[
 M=\sum_{x\in V}t(x),\quad
 P_k(x_1,\ldots,x_k)=\sum_{i=1}^{k-1}\langle x_i,x_{i+1}\rangle,
 \quad N_k=\sum_{P_k=0}\prod_{i=1}^k t(x_i).
\]
Then $QN_k=M^k+R_k$, where
\[
 R_k=\sum_{h\in\F_Q^\times}\sum_{x_1,\ldots,x_k\in V}
       \prod_i t(x_i)\,\psi(hP_k),
\]
and
\begin{equation}\label{eq:parseval}
 |R_k|\le (Q-1)Q^{d(k-1)/2}\|t\|_2^2\|t\|_\infty^{k-2}.
\end{equation}
No nonnegativity assumption on $t$ is required.
\end{finding}
\begin{proof}
Orthogonality gives $\sum_h\psi(hz)=Q\,1_{z=0}$; separating $h=0$ gives the identity. For $h\ne0$, set $K_h(x,y)=\psi(h\langle x,y\rangle)$ and $D=\operatorname{diag}(t)$. Nondegeneracy gives
\[
 (K_hK_h^*)_{x,x'}=\sum_y\psi(h\langle x-x',y\rangle)
 =Q^d1_{x=x'},
\]
so $\|K_h\|=Q^{d/2}$. The inner sum is
$t^{\mathsf T}K_hD K_h\cdots D K_ht$, with $k-1$ kernels and $k-2$ diagonal factors. The norm of the row functional $t^{\mathsf T}$ is $\|t\|_2$, even for complex $t$ (it is the Hermitian dual of $\bar t$). Cauchy--Schwarz and submultiplicativity give the bound for one $h$; sum the $Q-1$ terms. For $k=2$ there are no diagonal factors.
\end{proof}
This refines Fable suggestion~3's scalar toy in the exact form needed here. It is an open-path matrix coefficient, not in general a trace moment: tracing the power closes the path into a cycle.

\begin{finding}[\fid{F02}: finite spectral no-cancellation]\label{f:spectral}
For distinct nonzero complex numbers $\alpha_j$ and nonzero constants $c_j$, let $p_n=\sum_jc_j\alpha_j^n$. For a nonempty finite spectrum,
\[
 \limsup_{n\to\infty}|p_n|^{1/n}=\max_j|\alpha_j|.
\]
Consequently $|p_n|\le Cn^A R^n$ for all $n\ge1$, $R>0$, implies $|\alpha_j|\le R$. The same holds for any subexponential prefactor.
\end{finding}
\begin{proof}
Near zero,
\[
 \sum_{n\ge1}p_nz^n=\sum_j\frac{c_j\alpha_jz}{1-\alpha_jz}.
\]
Every $1/\alpha_j$ is a genuine pole, since the other summands are regular there. The radius of convergence is $1/\max|\alpha_j|$. Cauchy--Hadamard proves the equality; the bound gives radius at least $1/R$. The empty case is vacuous. Coincident spectral values must first be grouped; the conclusion concerns their \emph{nonzero} total coefficients.
\end{proof}
Positive multiplicities are one sufficient way to ensure that regrouping does not erase a root. Eigenvalues themselves need not be positive, or real. The traces of powers do not see nilpotent Jordan off-diagonal entries, so semisimplicity is not required for this lemma.

\begin{finding}[\fid{F03}: tensor-path purity criterion]\label{f:criterion}
Fix a prime power $q>1$, an integer $a$, $d\ge1$, and $B>0$. For every $n\ge1$ let $Q=q^n$ and $t_n:\F_Q^d\to\C$ satisfy $|t_n(x)|\le BQ^{a/2}$. Form $M_n,N_{k,n},R_{k,n}$ as in \fid{F01}. Assume:
\begin{enumerate}[label=(S\arabic*)]
\item A fixed finite multiset $\mathcal A$ of nonzero algebraic numbers, with positive integer multiplicities, and $\epsilon\in\{1,-1\}$ satisfy
\[
 M_n=\epsilon\sum_{\alpha\in\mathcal A}\alpha^n\qquad(n\ge1).
\]
\item For an unbounded set of integers $k\ge2$ there is, for each $k$, a fixed finite multiset $\mathcal B_k$ of nonzero algebraic numbers, possibly empty, with positive integer multiplicities such that
\[
 q^nN_{k,n}=-\epsilon^k\sum_{\beta\in\mathcal B_k}\beta^n\qquad(n\ge1).
\]
\end{enumerate}
Then $|\alpha|\le q^{(a+d)/2}$ for every $\alpha\in\mathcal A$.

For the all-conjugate assertion, require for every embedding $\sigma$ of the algebraic spectral data a corresponding trace system $t_{n,\sigma}$ satisfying these hypotheses with spectra $\sigma(\mathcal A),\sigma(\mathcal B_k)$ and a constant $B_\sigma$ independent of $n$. If in addition
\begin{equation}\label{eq:duality}
 \mathcal A=q^{a+d}/\mathcal A
 \quad\text{as multisets},
\end{equation}
then every eigenvalue has every complex absolute value $q^{(a+d)/2}$: the spectrum is pure of weight $a+d$.
\end{finding}
\begin{proof}
The pointwise bound gives $\|t_n\|_2^2\le B^2Q^{a+d}$. Equation~\eqref{eq:parseval} therefore gives
\begin{equation}\label{eq:overhead}
 |R_{k,n}|\le B^kQ^{k(a+d)/2+d/2+1}.
\end{equation}
By the exact path identity and (S1)--(S2),
\begin{equation}\label{eq:positivespectrum}
 -\epsilon^kR_{k,n}=
 \sum_{\beta\in\mathcal B_k}\beta^n+
 \sum_{(\alpha_1,\ldots,\alpha_k)\in\mathcal A^k}
 (\alpha_1\cdots\alpha_k)^n.
\end{equation}
After grouping equal roots all coefficients are positive integers. In particular the diagonal tuple contributes $\alpha^k$ and cannot be cancelled. For each \emph{fixed} $k$, \fid{F02} and \eqref{eq:overhead} give
\[
 |\alpha|^k\le q^{k(a+d)/2+d/2+1},\qquad
 |\alpha|\le q^{(a+d)/2+(d/2+1)/k}.
\]
Let $k$ tend to infinity through the unbounded set. This order of limits is essential: constants such as $B^k$ disappear in the $n$th-root limit before $k$ varies. Applying the same argument to each $\sigma$ gives all-conjugate upper bounds. Since $a+d$ is an integer, $q^{a+d}$ is rational and fixed by every $\sigma$. Apply the upper bound to $q^{a+d}/\sigma(\alpha)$, which occurs by \eqref{eq:duality}, to obtain the reverse inequality.
\end{proof}
The purity conclusion is conditional on a \emph{specified Tate duality}. An early draft allowed an arbitrary algebraic dualizing scalar; the independent check found that its conjugates could have different absolute values. The final theorem uses the integer-weight scalar $q^{a+d}$ and the correction is recorded in the audit.

\begin{finding}[\fid{F04}: controlled background terms]\label{f:errors}
Replace (S2) by
\[
 q^nN_{k,n}=-\epsilon^k\sum_{\beta\in\mathcal B_k}\beta^n+E_{k,n}.
\]
Suppose $|E_{k,n}|\le C_kn^{A_k}q^{n(k(a+d)/2+\eta_k)}$ for all $n$, with $\eta_k/k\to0$ along the unbounded set of $k$. Then \fid{F03}'s upper bound, and its purity conclusion under the same all-embedding duality hypotheses, still hold, provided that the modified identity and error estimate also hold for every corresponding trace system $t_{n,\sigma}$, with constants $C_{k,\sigma},A_{k,\sigma}$ and $\eta_{k,\sigma}/k\to0$ allowed to depend on $\sigma$.
\end{finding}
\begin{proof}
The positive spectral sum in \eqref{eq:positivespectrum} is now $\epsilon^k(E_{k,n}-R_{k,n})$. Its exponential rate for fixed $k$ is at most
$q^{k(a+d)/2+\max(d/2+1,\eta_k)}$. Apply \fid{F02} and take the $k$th root. Since $\max(d/2+1,\eta_k)/k\to0$, the same limit follows. Duality supplies the reverse bound as before.
\end{proof}
This is a sufficient error budget, not a necessity theorem for every possible proof. It permits suitable removed Tate or boundary terms; it does not justify discarding uncontrolled terms.

\begin{finding}[\fid{F17}: finite extensions and a separate dual system]\label{f:extensiondual}
The upper-weight conclusion of \fid{F03}--\fid{F04} still holds if, for each allowed $k$, the opposite-sign identity and its error estimate are known only for $n=e_kr$, all $r\ge1$, where the positive integer $e_k$ may depend on $k$.

Self-duality can also be replaced by a separate dual coefficient system. Suppose the first spectrum $\mathcal A$ has the proved upper weight $a+d$, a second spectrum $\mathcal A^D$ has the proved upper weight $d-a$, and a perfect Frobenius-equivariant pairing identifies
\[
 \mathcal A^D=\{q^d/\alpha:\alpha\in\mathcal A\}
\]
with multiplicities. Then $\mathcal A$ is pure of weight $a+d$, provided both upper bounds hold under every embedding.
\end{finding}
\begin{proof}
For fixed $k$, apply the spectral lemma in the variable $r$. The positive multiset now contains $\alpha^{e_kk}$; equal powered roots still have positive multiplicities, so none disappears. Its exponential bound is
\[
 (q^{e_k})^{k(a+d)/2+d/2+1}
\]
(or the corresponding maximum with $\eta_k$). Taking the $(e_kk)$th root gives exactly the same bound on $|\alpha|$ as before, and the unbounded-$k$ limit is unchanged. Constants may depend on $k$ and $e_k$. This uses full multiples of $e_k$, not an arbitrary infinite subsequence such as the odd degrees in \fid{F05}(d).

For the second assertion apply the dual upper bound to $q^d/\alpha$:
\[
 |\alpha|\ge\frac{q^d}{q^{(d-a)/2}}=q^{(d+a)/2}.
\]
Combine it with the first upper bound. Since $q^d$ is rational, the same calculation holds under every embedding.
\end{proof}
This refinement allows a geometric construction to choose a finite field of definition separately for each path length. The upper weight $d-a$ for the dual system is an independent hypothesis: the original pointwise upper weight $a$ alone does not give the required upper estimate for the dual coefficient system. The refinement also distinguishes Verdier duality with a dual coefficient system from an assertion that a single compact-support module is self-dual. To use the latter option geometrically, the necessary compact-support/ordinary-cohomology comparison still has to be proved; it is not supplied by dimension alone.

\subsection{What the hypotheses mean geometrically}\label{sec:tracegeom}
Let $\mathcal F$ be a constructible Weil sheaf on $\A^d/\F_q$. Its trace functions $t_n$ have the required pointwise estimate if stalk ranks are bounded by $B$ and all stalk eigenvalues have weight at most $a$. Let
\[
 H_k=\{(x_1,\ldots,x_k):P_k(x_1,\ldots,x_k)=0\}
 \subset(\A^d)^k.
\]
The character kernel is represented by an Artin--Schreier sheaf pulled back along $hP_k$; the coefficient system on $H_k$ is the restriction of $\mathcal F^{\boxt k}$. The Grothendieck trace formula identifies the sums with alternating traces of $R\Gamma_c$.

If $R\Gamma_c(\A^d_{\bar{\F}_q},\mathcal F)$ is concentrated in degree $d$, then $\epsilon=(-1)^d$. If the relevant cohomology of $H_k$ is concentrated in degree $kd-1$, its sign is
\[
 (-1)^{kd-1}=-\bigl((-1)^d\bigr)^k.
\]
Multiplication by $q^n$ is the twist $(-1)$. These facts give precisely (S1)--(S2), or \fid{F04} when background terms have been controlled. A suitable duality of the coefficient system and cohomology supplies \eqref{eq:duality}.

Being a hypersurface of dimension $kd-1$ does \emph{not} prove this concentration. Singularities, exceptional fibers, compact-support boundary groups and Tate contributions are separate obligations. An affine variety also need not identify $H_c^d$ with its dual without an additional comparison. This is the information the proposed sign route still needs. Its spectral core can be supplied geometrically or by a rigorous algebraic/category substitute; our argument does not prove that schemes are the only possible source of it.

\subsection{Countermodels to omitted hypotheses}
\begin{finding}[\fid{F05}: precise failures of abbreviated sign arguments]\label{f:countermodels}
The following general implications are false.
\begin{enumerate}[label=(\alph*)]
\item Path identity plus Parseval implies the expected middle weight.
\item Bounds for alternating traces control every individual cohomological eigenvalue without non-cancellation.
\item An upper-weight bound implies purity of that weight.
\item Bounds on an arbitrary infinite set of extension degrees suffice.
\item Nonnegative weights and Parseval make the Fourier operator positive.
\item Punctual purity implies arithmetic Frobenius semisimplicity.
\end{enumerate}
\end{finding}
\begin{proof}
(a) Take $t_n=1$ on $\F_{q^n}$, $a=0,d=1$. Then $M_n=q^n$, with eigenvalue $q>\sqrt q$, while \fid{F01} holds for all $k,n$. The actual cohomology is $H_c^2(\A^1)=\Q_\ell(-1)$, not the assumed middle degree.

(b) On $\Spec\F_q$, the complex $\Q_\ell(-1)\oplus\Q_\ell(-1)[-1]$ has a $q$-eigenvalue in degrees 0 and 1. All alternating traces vanish, and its virtual zeta function is 1.

(c) The eigenvalue 1 satisfies the weight-at-most-1 bound but has actual weight zero.

(d) The spectrum $\{R,-R\}$ has trace zero for every odd $n$, with $R$ arbitrarily large. Our lemma uses every extension degree; we assert no general necessity theorem excluding other suitably controlled subsequences.

(e) For $t=1$ on $\F_3$, the normalized Fourier matrix has diagonal entry $e^{2\pi i/3}/\sqrt3$ at $x=1$ and hence is not Hermitian. Over $\F_2$ the normalized matrix is $2^{-1/2}\left(\begin{smallmatrix}1&1\\1&-1\end{smallmatrix}\right)$, with eigenvalues $1,-1$. Both preserve the Euclidean norm.

(f) The matrix $\left(\begin{smallmatrix}1&1\\0&1\end{smallmatrix}\right)$ has pure weight-zero eigenvalues and trace 2 for every power, but is not diagonalizable. It defines a continuous $\ell$-adic representation through $\widehat\Z\to\Z_\ell$ and the associated unipotent matrices. This is the standard arithmetic-versus-geometric distinction also discussed in BBD after 5.3.8.
\end{proof}
In particular Fable suggestion~3's proposed identification of Fourier-after-multiplication with a positive operator is false in that literal generality. This does not refute the valid norm estimate or an independently constructed different positive operator.

\section{From purity to positivity: what is additionally required}\label{sec:positivity}
\begin{finding}[\fid{F06}: weights, positive values, and unimodality are different]\label{f:phase}
Exact purity does not imply a positive Frobenius trace. Even positivity at every positive numerical specialization, together with bar invariance, does not imply coefficientwise positivity. Pure graded pieces with a duality need not have unimodal dimensions.
\end{finding}
\begin{proof}
On $\Spec\F_q$, the rank-one Weil sheaf with Frobenius $-1$ is pure of weight zero but has trace $(-1)^n$ over $\F_{q^n}$. Next
\[
 h(v)=v^2-1+v^{-2}\ge1\quad(v\in\mathbb R\setminus\{0\})
\]
by $v^2+v^{-2}\ge2$, but its central coefficient is negative and $h(v^{-1})=h(v)$. Since $h(1)=1$, a factorization $\pm(v-v^{-1})^a f(v)$ with $a\ge0$ and nonnegative Laurent $f$ would force $a=0$, positive sign and $f=h$, which is impossible. Finally take $M_0=\Q_\ell$, $M_1=0$, $M_2=\Q_\ell(-2)$, with grade $j$ assigned weight $2j$ and the perfect pairing $M_0\otimes M_2\to\Q_\ell(-2)$. Its dimension sequence is $(1,0,1)$, not unimodal. Every degree-one operator $L$ has $L^2=0$, so it cannot satisfy the isomorphisms of \fid{F07} for $D=2$. This fixes the grading convention; it is not the usual degree-two Lefschetz grading on the cohomology of $\PP^1$.
\end{proof}
The last example is an abstract graded Weil object, not the cohomology of a smooth projective variety to which Hard Lefschetz already applies. It isolates precisely the extra theorem that such geometry would supply.

\begin{finding}[\fid{F07}: effective categorification and Lefschetz suffice]\label{f:category}
Let $\mathscr C$ be a split semisimple additive category with grading shift $\langle1\rangle$ and an additive product. Suppose its designated simple representatives satisfy finite decompositions
\[
 L_a*L_b\simeq\bigoplus_{c,j}L_c\langle j\rangle^{\oplus m_{ab}^{c,j}}.
\]
If an isomorphism from its graded split Grothendieck ring to the target algebra sends $[L_a]$ to the \emph{specified} basis element $b_a$ and grading shift to $v$, then its structure constants are
$\sum_jm_{ab}^{c,j}v^j\in\N[v,v^{-1}]$.

For a normalized multiplicity space $M=\bigoplus_{j=0}^D M_j$, suppose there is an operator $L:M_j\to M_{j+1}$ such that
$L^{D-2j}:M_j\to M_{D-j}$ is an isomorphism for $j\le D/2$. Then $\sum_j\dim M_jv^j$ is symmetric and unimodal.
\end{finding}
\begin{proof}
Apply the split Grothendieck relation to the displayed direct sum and compare coefficients under the basis-preserving isomorphism. Multiplicities of actual direct summands are nonnegative integers. For the second assertion the displayed isomorphisms give symmetry. If $j<D/2$ and $Lx=0$, then $L^{D-2j}x=0$, so $x=0$. Thus $\dim M_j\le\dim M_{j+1}$ up to the middle; symmetry gives the other half.
\end{proof}
For VIII this must be applied to the correctly \emph{normalized} multiplicity polynomial, with the signs, $(v-v^{-1})$ factors and positive rational denominators justified separately. The grading step must match the variable in which unimodality is asserted. An arbitrary derived $K_0$ realization is insufficient because $[K[1]]=-[K]$.

\begin{finding}[\fid{F08}: pure geometry and positive Gram determinants do not fix a basis]\label{f:twopoints}
There is a smooth projective finite \(\F_q\)-scheme, an integral basis of effective $K_0$ classes, and a positive definite Gram form, such that multiplication has a negative structure constant in that basis.
\end{finding}
\begin{proof}
Take $X=\Spec\F_q\sqcup\Spec\F_q$. Its $K_0$ ring is $\Z^2$ with componentwise product. The realization $H^0_{\et}(X_{\bar{\F}_q},\Q_\ell)=\Q_\ell^2$ has identity Frobenius and is pure of weight zero. Put
\[
 b_1=(1,1),\qquad b_2=(1,2).
\]
They form an integral basis, since the determinant is 1, and both are actual vector bundles, with the displayed ranks on the two components. The natural real pairing has Gram matrix
\[
 \begin{pmatrix}2&3\\3&5\end{pmatrix},
\]
whose principal minors are $2,5,1>0$. Nevertheless
\[
 b_2^2=(1,4)=3b_2-2b_1.
\]
All assertions follow directly. Exterior powers are also defined componentwise: $\lambda_t(r,s)=((1+t)^r,(1+t)^s)$. Thus passing to the ordinary geometric $\lambda$-ring framework does not remove this counterexample.
\end{proof}
This disproves a general bridge from purity/determinants/effectivity to coefficients of a chosen basis. It does not disprove the positivity of any designated GCT canonical basis.

\section{An obstruction in the actual nonstandard Hecke inclusion}\label{sec:ambient}
\begin{finding}[\fid{F09}: the inherited positive cone has rank one]\label{f:hecke}
Let $R=\Z[v,v^{-1}]$, $R_+=\N[v,v^{-1}]$, $\delta=v+v^{-1}$, and let $H=R1\oplus Rc$ with $c^2=\delta c$. In $H\otimes_RH$, put
\[
 P=c\otimes c+(c-\delta)\otimes(c-\delta),\qquad A=R[P].
\]
Equip the ambient tensor algebra with its positive cone $C_+$ in the basis
$1,c\otimes1,1\otimes c,c\otimes c$. Then
\begin{equation}\label{eq:cone}
 A=R1\oplus RP,\qquad P^2=\delta^2P,\qquad A\cap C_+=R_+1.
\end{equation}
Consequently no basis of the rank-two algebra $A$ consists entirely of ambient-positive elements, although its internal basis $\{1,P\}$ has positive multiplication.
\end{finding}
\begin{proof}
Set $b=c-\delta$. Then $cb=0$ and $b^2=-\delta b$, so the cross terms vanish and
$P^2=\delta^2(c\otimes c+b\otimes b)=\delta^2P$. Also
\[
 P=\delta^2\,1-\delta(c\otimes1+1\otimes c)+2c\otimes c.
\]
This proves that $1,P$ span $A$ and are independent. For $a+bP$, its four ambient coefficients are
\[
 a+\delta^2b,\quad-\delta b,\quad-\delta b,\quad2b.
\]
If all lie in $R_+$, then $2b\in R_+$ implies coefficientwise $b\ge0$. A nonzero such $b$ makes $\delta b$ nonzero and nonnegative, contradicting $-\delta b\ge0$. Thus $b=0$, and then $a\in R_+$. The converse is immediate. Finally $\delta^2=v^2+2+v^{-2}$ is positive, proving the intrinsic positivity assertion.
\end{proof}
This is the actual rank-one nonstandard generator of revised IV, \S1.4, equation~(8), and Definition~11.1, equation~(93). We fix the consistent earlier \S3.1 convention
$(T-v)(T+v^{-1})=0$, $c=C'_s=T+v^{-1}$ and $C_s=T-v=c-\delta$.
The opening of \S11.1 interchanges the exponents while retaining the idempotent assertion; this printed normalization slip is recorded rather than silently used. The proof above is unambiguous and also works over $\Q[v,v^{-1}]$ with the rational positive cone.

\paragraph{Categorical consequence.}
Suppose the usual graded tensor Hecke category has split $K_0$ identified with $H\otimes H$, with actual objects giving the ambient positive cone. Every object in an additive graded subcategory whose character lies in $A$ has class in $A\cap C_+=R_+1$. Its Grothendieck image therefore spans at most $R1$, not $A$. Hence the full nonstandard algebra cannot be obtained as such an effective subcategory through this fixed inclusion.

This is a proved obstruction to a specific natural route. It allows a different category, a different heart or character map, or a derived realization with signs. It is not a counterexample to VIII's signed-normalized conjecture and does not assert that $P$ must be a designated canonical-basis vector. The contrast between intrinsic positivity and failure of inherited positivity is exactly why a new realization is needed for this approach.

\begin{finding}[\fid{F10}: an algebraic replacement and its limits]\label{f:pivots}
Let $R=\Q[v,v^{-1}]$ and $R_+=\Q_{\ge0}[v,v^{-1}]$. Suppose an $R$-algebra $B$ has a positive basis $(e_j)$ and an $R$-subalgebra $A$ has a basis $(b_i)$ with expansions $b_i=\sum_jp_{ji}e_j$, $p_{ji}\in R_+$. If each $b_i$ has an exclusive coordinate $j_i$ with coefficient $r_iv^{m_i}$, $r_i>0$, and all other $b_h$ have zero coefficient there, then $A$ has positive structure constants in this basis. Disjoint positive supports, without this pivot normalization, do not suffice, even with bar invariance and generic semisimplicity. Disjoint supports only after specialization are weaker still.
\end{finding}
\begin{proof}
Write $b_sb_t=\sum_i m_{st}^i b_i$. Its ambient coefficients are positive. The exclusive coordinate is $r_iv^{m_i}m_{st}^i$, so division by the positive monomial preserves the cone and proves $m_{st}^i\in R_+$.

For the second assertion put
\[
 p=v+v^{-1},\quad h=v^2-1+v^{-2},\quad c=ph=v^3+v^{-3},
\]
and take $B=R\times R[t]/(t^2-c)$ with basis
$1=(1,1)$, $x=(0,t)$, $y=(0,1)$. It has positive multiplication
$x^2=cy$, $xy=x$, $y^2=y$. The free subalgebra with basis $1,x,z=py$ has disjoint positive ambient supports but
\[
 x^2=hz,\quad xz=px,\quad z^2=pz.
\]
The coefficient $h$ is negative in its middle term, and \fid{F06} shows it also fails the stated signed-factor positivity form. The product presentation proves associativity, and $t^2-c$ is separable over $\Q(v)$, so the generic algebra is semisimple. Since $p$ is invertible over $\Q(v)$, the generic subalgebra equals the generic ambient algebra and is also semisimple. All the displayed data are bar-invariant.

Finally in the positive algebra $\Q[v][x]/(x^2-x)$, the basis $b_0=1$, $b_1=x+v$ has nonnegative external coefficients and disjoint singleton supports modulo $v$, but direct expansion gives
\[
 b_1^2=(1+2v)b_1-v(1+v)b_0.
\]
This proves the last assertion.
\end{proof}
The statement provides a possible algebraic route that does not require new finite-field geometry. VIII, Conjecture~2.2, asserts disjoint supports in $L_\alpha/vL_\alpha$, not the global exclusive monomial pivots of \fid{F10}. Nor is raw positivity necessarily appropriate for all of VIII's signed targets. The missing pivot/normalization theorem cannot be obtained merely by repeating the local support assertion.

\begin{finding}[\fid{F11}: two limits of the Grothendieck-ring shortcut]\label{f:lambda}
Neither virtual $\lambda$-ring identities nor reduction of a noncommutative algebra modulo $p$ supplies the missing effective positive cone and arithmetic Frobenius.
\end{finding}
\begin{proof}
In $K_0(\Rep\mathbb G_m)=\Z[z,z^{-1}]$, the line character $z$ satisfies
\[
 \lambda_t(z-1)=\frac{1+zt}{1+t}
 =1+(z-1)t+(1-z)t^2+\cdots.
\]
Thus $\lambda^2(z-1)=1-z$, outside the effective cone $\N[z,z^{-1}]$. Also
$\gamma_t(z-1)=\lambda_{t/(1-t)}(z-1)=1+(z-1)t$. The augmentation/\(\gamma\)-filtration bookkeeping supplies no additional positivity order.

For the Frobenius assertion take $A=e_{12}$ and $B=e_{21}$ in $M_2(\F_p)$. Their $p$th powers vanish, whereas $(A+B)^2=I$, hence $(A+B)^p=I$ for $p=2$ and $A+B$ for odd $p$. The map $x\mapsto x^p$ is not additive. Therefore the automatic absolute Frobenius of a commutative scheme is not obtained by the same formula on a noncommutative coordinate algebra.
\end{proof}
A separately constructed quantum or categorical Frobenius remains possible. Also braided tensor categories do not automatically carry the symmetric exterior-power structure required by SGA~6.

\section{The theta route: exact positive consequences}\label{sec:theta}
\subsection{Periods and multiplicities are different data}
Alp\"oge's step~4 (p.~3) asserts positivity of a functional via a determinant. Step~7 (p.~4) describes determinants containing GGP periods, interlacing values 0 or 1, and controlled failures of interlacing. Step~8 describes iterated lifts between totally definite unitary groups and a gain of multiplicity. Page~1 and the introductory paragraph on p.~3 explain this as the analogue of Stepanov order of vanishing, not an assertion about the dimension of a Kronecker Hom space. The supplied Lean wrappers contain no theorem identifying the two kinds of multiplicity.

The following two findings state exactly what a positive projector or a chain of branching rules can give. They also show why merely finding a positive determinant is not the missing representation-theoretic formula.

\begin{finding}[\fid{F12}: positive projectors already realize numerical multiplicities]\label{f:projector}
For a finite group $H$ and a finite-dimensional complex unitary representation $U$, put
\[
 P_H=\frac1{|H|}\sum_{h\in H}\rho(h).
\]
It is the orthogonal projector onto $U^H$, with $\Tr P_H=\rank P_H=\dim U^H$ and nonnegative principal minors in an orthonormal basis. For partitions of $D$,
\[
 g(\lambda,\mu,\nu)=\Tr P_{S_D}
 \quad\text{on }[\lambda]\otimes[\mu]\otimes[\nu].
\]
For $\lambda\vdash dn$ and $\dim V\ge\ell(\lambda)$,
\[
 a_\lambda(d[n])=\Tr P_{H_{d,n}}\quad\text{on }[\lambda],\qquad
 H_{d,n}=(S_n)^d\rtimes S_d\subset S_{dn},
\]
where $a_\lambda(d[n])$ is the multiplicity in $\Sym^d(\Sym^nV)$.
\end{finding}
\begin{proof}
Inversion gives $P_H^*=P_H$, and counting each element $|H|$ times in the double sum gives $P_H^2=P_H$. Its image is fixed by $H$, and it is the identity on invariant vectors. Thus it is the claimed orthogonal projector. In an orthonormal basis its principal submatrices are Gram matrices of the vectors $P_He_i$, so their determinants are squared exterior norms. Its spectrum consists of $\dim U^H$ ones and the remaining zeros.

Specht modules are self-dual and can be made unitary by averaging a form; the triple invariant-space description gives the Kronecker assertion. For plethysm, invariants under the within-block permutations give $(\Sym^n V)^{\otimes d}$, and under block permutations give $\Sym^d(\Sym^nV)$. Apply invariants to Schur--Weyl decomposition
\[
 V^{\otimes dn}=\bigoplus_\lambda S_\lambda V\otimes[\lambda]
\]
and compare the $S_\lambda V$ multiplicity. This is also the projector explicitly used in BIP, equation~(4.1), with Proposition~3.3.
\end{proof}
This is a proof of the useful reformulation, not a new numerical nonnegativity result. BIP, Proposition~4.5, already gives projected tableaux spanning the highest-weight space; its Theorem~4.7 expresses evaluations as sums of signed column contributions. Independence and a compatible basis remain separate.

For any spanning vectors $v_i$ of a multiplicity space, their Gram matrix satisfies
\[
 \sum_{|I|=r}\det G_I=\sum_{|I|=r}\bigl\|\bigwedge_{i\in I}v_i\bigr\|^2.
\]
This sum is positive exactly when the rank is at least $r$. Determinant positivity therefore restates the rank question; it does not by itself determine which minors are nonzero or provide an efficient basis-selection rule. We make no claim that every such rank question is intrinsically hard.

\begin{finding}[\fid{F13}: when an interlacing chain counts a target multiplicity]\label{f:branching}
Let $\mathscr C_r,\ldots,\mathscr C_0$ be split semisimple linear categories of finite-length objects and let $F_i:\mathscr C_i\to\mathscr C_{i-1}$ be additive exact functors. Suppose
\[
 F_i S_i(a)\simeq\bigoplus_b S_{i-1}(b)^{\oplus b_i(a,b)},\qquad b_i(a,b)\in\{0,1\},
\]
with finite sums. Then the multiplicity of $S_0(c)$ in $F_1\cdots F_rS_r(a)$ is the number of allowed paths from $a$ to $c$:
\[
 \sum_{a_{r-1},\ldots,a_1}
 b_r(a,a_{r-1})\cdots b_1(a_1,c).
\]
If the allowed steps are specified by integral interlacing inequalities and every intermediate coordinate is bounded by the fixed endpoints, these paths are the integer points of the resulting polytope.
\end{finding}
\begin{proof}
Apply each additive functor to the previous finite direct-sum decomposition. The multiplicities multiply along paths and add over intermediate simples. Semisimplicity identifies the final coefficient with the desired multiplicity. Every product of 0/1 entries is 1 exactly for an allowed path. Sending a path to its array of intermediate integer weights is a bijection with the integer solutions of the inequalities. Boundedness makes the polyhedron a polytope.
\end{proof}
Ordinary restriction along consecutive block inclusions $\GL_r\supset\GL_{r-1}$ has the interlacing rule $\lambda_j\ge\mu_j\ge\lambda_{j+1}$; see \cite[Theorem 2.1, equation (2.2), and the patterns (2.4)]{Molev}. The Kronecker map $(g,h)\mapsto g\otimes h$ and the plethysm map $g\mapsto S_\mu(g)$ are different maps. A theta lift between dual pairs is not automatically their composite restriction functor.

A small object mismatch is instructive. $S_{(d^r)}(\C^r)=\det^d$ is one-dimensional and has a unique usual interlacing pattern, whereas the rectangular Specht module is a different representation. For example
\[
 [(2,2)]\otimes[(2,2)]=[(4)]\oplus[(2,2)]\oplus[(1^4)].
\]
To check this without a table lookup, the permutation action of $S_4$ on the three pairings of four letters is the direct sum of the trivial representation and $[(2,2)]$. The latter has character $(2,0,2,-1,0)$ on cycle types $(1^4),(2,1,1),(2,2),(3,1),(4)$. Its pointwise square equals this character plus the trivial and sign characters, proving the identity. A correctly identified chain could still count higher multiplicities; the example only prevents substitution of a determinant character for a rectangular Specht module.

\subsection{A lattice obstruction to ordinary interlacing alone}
\begin{finding}[\fid{F18}: homogeneous difference constraints and a rectangular parity ray]\label{f:saturation}
Let $u$ range over an integral label set closed under positive-integer scaling. Fix a number of coordinates and a finite list of inequalities
\[
 P(u)=\{x\in\mathbb R^s:x_i-x_j\le b_{ij}(u)\},\qquad x_0=0,
\]
where every specified $b_{ij}(u)$ is integral and $b_{ij}(Nu)=Nb_{ij}(u)$ for every $N\ge1$. Assume $P(u)$ is bounded, and put $h(u)=\#(P(u)\cap\Z^s)$. Then
\[
 h(Nu)>0\quad\Longleftrightarrow\quad h(u)>0.
\]
The same property holds for a fixed finite sum of such counting functions. Thus these models cannot count all two-row rectangular Kronecker coefficients. On the particular ray with all three partitions $(d,d)$, the established two-row formula gives
\[
 g((d,d),(d,d),(d,d))=\boldsymbol1_{2\mid d}
 =\#\bigl(d\{1/2\}\cap\Z\bigr)\qquad(d\ge1).
\]
This ray has an exact rational Ehrhart model despite its failure of the displayed ordinary support property.
\end{finding}
\begin{proof}
Homogeneity gives $P(Nu)=NP(u)$. If $x_i-x_j\le b$ with $b\in\Z$, then
$\lfloor x_i\rfloor\le\lfloor x_j+b\rfloor=\lfloor x_j\rfloor+b$.
Hence coordinatewise flooring, including $x_0=0$, takes every real point of $P(u)$ to an integral one. If $z\in P(Nu)\cap\Z^s$, apply this to $z/N$. Conversely an integral $x\in P(u)$ gives $Nx\in P(Nu)\cap\Z^s$. Positivity of a fixed finite sum is positivity of at least one summand, so the same argument applies. Boundedness is only needed to make these finite counts.

For $u=((1,1),(1,1),(1,1))$, the triple tensor product is the sign representation of $S_2$, so $g(u)=0$. The explicit $S_4$ character calculation following \fid{F13} gives $g(2u)=1$: the class sizes are $(1,6,3,8,6)$ and the character $(2,0,2,-1,0)$ has cubic average $(8+3\cdot8-8)/24=1$. Thus an exact model of the stated form on all two-row rectangles is impossible.

For the whole ray use IV, \S18.3, equation~(172), printed p.~128. In its column convention, $(d,d)=[d,0]$ as a two-row label and $[0,0,d,0]$ as a four-row label. The coefficient for the second rectangle is the constant term of that equation's generating function. Only $k=0$ and $j=0$ in its first sum survive; the second sum is empty. The remaining condition is that $n_3,n_2,n_1$ are all even and nonnegative. Substituting $(n_3,n_2,n_1)=(0,d,0)$ gives exactly the displayed parity formula. The single rational point $\{1/2\}$ has an integral point after dilation by $d$ exactly when $d$ is even.
\end{proof}
The difference-constraint result includes ordinary interlacing with integral endpoint rows. Its exclusions are essential: extra congruences, sum equations, nonhomogeneous endpoint shifts, a changing constraint list or a changed lattice can evade it. For example $2z=d$ retains precisely the parity that ordinary integral differences lose. The example concerns all two-row rectangles, not a model restricted only to IP's large-parameter region.

Nor does it contradict VI's definition of ``strictly saturated'' (Definition~1.2.4, printed p.~14). That definition tests each nonzero residue polynomial separately; here $f_{\rm even}=1$ and $f_{\rm odd}=0$, so its saturation index is zero. We use the displayed support property to avoid conflating the two notions. IV, Theorem~18.2 and equations~(163)--(164), explicitly retain congruences and remove columns in pairs. A successful theta comparison must preserve these arithmetic data or equivalent information; the source does not identify them with a particular GGP local sign.

In Grothendieck's language, the tensor group-scheme map $\GL_{2,\Z}\times\GL_{2,\Z}\to\GL_{4,\Z}$ induces, over $\Q$, a positive restriction map on representation $K_0$. The coefficient of $[\det^d\boxtimes\det^d]$ in the restriction of $[S_{(d,d)}(V\otimes W)]$ is this parity function. Thus an effective $K_0$ homomorphism does not by itself give the ordinary support property. The proof above is already the exact integral-lattice dictionary; its group-scheme framework is SGA~3 I, 4.7.2--3, with cote~90, scan~17, as an action/comodule anchor.

\paragraph{Exact remaining theta gate.}
Construct functors and endpoint identifications making Alp\"oge's lift chain the target restriction, or identify its period spaces with the projector images of \fid{F12}. To obtain canonical-basis positivity, additionally identify grading, operations, and the specified basis as in \fid{F07}. For a \(\#P\) rule the path encodings and admissibility tests must have polynomial bounds; for PH1 the stronger polyhedral specifications in VI must also hold. None of these identifications is supplied by steps~4, 7 or~8.

\section{Two further geometric routes and their limits}\label{sec:other}
\begin{finding}[\fid{F14}: what an ambient pure realization would prove]\label{f:sub}
A Frobenius-stable subquotient of a pure algebraic Weil module of weight $w$ is pure of weight $w$. A stable subquotient of its endomorphism module, with Frobenius acting by conjugation, is pure of weight zero.
\end{finding}
\begin{proof}
A basis adapted to a stable subspace makes Frobenius block upper triangular. The characteristic polynomial factors into those of the subspace and quotient, so their eigenvalues are selected from the original spectrum. This works under every embedding. Eigenvalues on $V\otimes V^*$ are ratios $\alpha/\beta$, of absolute value 1; apply the first assertion. No diagonalizability is required.
\end{proof}
Thus if the desired nonstandard multiplicity objects were already realized as such stable subquotients, a new spectral RH would not be needed for them. Stability, grading, effective classes, and a basis-preserving realization would still need proof. Declaring an arbitrary vector space to have identity Frobenius proves an artificial weight-zero assertion and identifies no GCT coefficient.

\begin{finding}[\fid{F15}: orbit models exist, but their ordinary purity does not determine Hilbert multiplicities]\label{f:orbit}
An integral polynomial tensor $f$ in a finite free representation $W_\Z$ of $\GL_{N,\Z}$ has a finite-type affine orbit-closure model over $\Z$ whose complex fiber is the usual reduced orbit closure. Ordinary pure cohomology of a projective scheme does not determine its polarized coherent Hilbert function or its equivariant coherent character.
\end{finding}
\begin{proof}
The orbit map induces $\Z[W]\to\Z[\GL_N]$. Its kernel $I$ defines the finite-type scheme $\Spec(\Z[W]/I)$. The quotient is a subring of $\Z[\GL_N]$ and is torsion-free, hence flat over $\Z$. Since $\C$ is flat over $\Z$, tensoring preserves the kernel. The resulting image ring is a subring of the reduced domain $\C[\GL_N]$, so the complex fiber is the reduced scheme-theoretic orbit closure. No corresponding equality after arbitrary reduction modulo $p$ is asserted; kernel formation can fail to commute with nonflat base change. For an explicit example, let $\mathbb G_m$ scale $f=p$ in $\A^1_\Z$. The map $\Z[x]\to\Z[t,t^{-1}]$, $x\mapsto pt$, is injective, so this model is $\A^1_\Z$. Reducing $f$ first gives zero and an orbit closure consisting of one point, whereas the special fiber of the model is $\A^1_{\F_p}$.

For the second assertion fix $X=\PP^1_{\F_q}$ and $L_r=\mathcal O(r)$, $r>0$. Monomials of degree $rd$ in two variables form a basis of $H^0(X,L_r^d)$, so its dimension is $rd+1$. The underlying scheme and its ordinary/intersection cohomology are independent of $r$: they are $\Q_\ell$ in degree 0 and $\Q_\ell(-1)$ in degree 2 (the projective-space calculation, \cite[Example 16.3, p.~109]{Milne}). Thus the pure Frobenius data do not determine the polarized Hilbert function. Finally, on a point with trivial $\mathbb G_m$ action, a weight-$r$ linearization of $\mathcal O$ gives character $z^{rd}$ on its $d$th tensor power, while the underlying Weil cohomology remains the same. This proves the character assertion.
\end{proof}
A projective integral model of the homogeneous orbit closure can likewise be formed by projectivization. For a reduced equidimensional projective fiber $X/\F_q$ of dimension $e$ with a dense smooth open $j:U\hookrightarrow X$, the unshifted complex $\IC_X^u=j_{!*}\Q_\ell[e][-e]$ is pure of weight zero (BBD~5.3.2), and $IH^i(X)=H^i(X_{\bar{\F}_q},\IC_X^u)$ is pure of weight $i$ by proper pushforward and duality (BBD, Remark~5.4.9, printed p.~143). For a nonreduced fiber pass to its reduction for this \(\ell\)-adic assertion. This is already ordinary finite-field mathematics. The GCT multiplicities instead belong to the graded coherent representation $\C[X]_d$, or to coherent sections with the required saturation and degree hypotheses. The polarization, linearization and comparison object cannot be omitted. Riemann--Roch is consistent with the example: it explicitly retains the varying Chern character of $L_r^d$.

Hilbert/Quot schemes represent the relevant coherent moduli under their usual hypotheses (FGA, Expos\'e~221, Theorems~3.1--3.2); they do not identify their intersection cohomology with the desired multiplicity spaces. Constructing an affine space from an already known multiplicity module encodes its unknown dimension in a scheme, but does not compute that dimension or produce a canonical indexing set. This route was pursued and reduces to an additional coherent-to-constructible comparison, not to an unproved ordinary Weil theorem.

\begin{finding}[\fid{F16}: a positive rank-one model exists with a different action]\label{f:rankonemodel}
The abstract rank-one algebra of \fid{F09} has a positive realization inside the same ambient algebra after changing the inclusion. Put $E=c\otimes c$. Then $R1\oplus RE$ is a subalgebra, and $P\mapsto E$ gives an $R$-algebra isomorphism from $R1\oplus RP$. The coefficient $\delta^2$ is the normalized Tate cohomology polynomial of $\PP^1\times\PP^1$. This isomorphism does not preserve the original ambient action relevant to GCT.
\end{finding}
\begin{proof}
We have $E^2=\delta^2E$, and $1,E$ are independent in the ambient tensor basis. Hence both algebras have the presentation $R[z]/(z^2-\delta^2z)$, with the respective chosen generators. The new basis $1,E$ is ambient-positive. The even cohomology of $\PP^1\times\PP^1$ is $\Q_\ell$, $\Q_\ell(-1)^{\oplus2}$ and $\Q_\ell(-2)$ in degrees 0, 2 and 4, by the projective-space calculation and K\"unneth. Its polynomial, normalized by $v^{-2}$ with $v^2$ counting a Tate degree, is
\[
 v^{-2}(1+2v^2+v^4)=v^{-2}+2+v^2=\delta^2.
\]
At $v=1$, let $s$ be the nontrivial element of $S_2$, so $c=1+s$. Then
\[
 P=2(1+s\otimes s),\qquad E=(1+s)\otimes(1+s).
\]
On the sign tensor sign representation, $P$ acts by 4 and $E$ by 0. Therefore the isomorphism changes the given ambient representation and cannot be substituted into the Kronecker comparison without a new compatibility theorem.
\end{proof}
This is a positive result for the abstract rank-one Hecke case and a precise limit of that result. Purity here is degree by degree: the total ungraded cohomology has several weights, not one weight-zero Weil module. The Tate-grade dimensions are $(1,2,1)$, whereas consecutive integer powers of $v$ give $(1,0,2,0,1)$; no assertion about VIII's precise unimodality normalization follows without fixing the grading step. This is neither an identification of VIII's full specified bases nor a new consequence of the sign-only theorem. In motive language the coefficient is realized by the Tate object $\mathbf1\oplus\mathbb L^{\oplus2}\oplus\mathbb L^2$, with a grading shift; the missing datum is the action comparison, not its purity.

\subsection{The actual rank-three algebra has a coherent scheme model}
\begin{finding}[\fid{F19}: a finite-flat center, its collision, and the lost radical]\label{f:rankthree}
Put $\delta=v+v^{-1}$, $D=\delta^2$, and take either
\[
 R=\Q[v,v^{-1}]_{(v-1)}
 \quad\text{or}\quad
 R_0=\Z[1/2,v,v^{-1},\delta^{-1},(D-1)^{-1},(D-2)^{-1}].
\]
Write $R$ for the chosen ring in the formulas below. Let $A$ be the actual nonstandard Hecke algebra $\check H_3$ over this ring, with its specified generators $P_1,P_2$. Define
\[
 z_1=\frac{D-1}{D},\qquad z_2=\frac{4(D-1)}{D^2},\qquad
 C=R[z]/((z-z_1)(z-z_2)).
\]
Then there is an explicit algebra isomorphism
\begin{equation}\label{eq:r3model}
 A\simeq R\times R\times M_2(C).
\end{equation}
On the two scalar blocks $(P_1/D,P_2/D)$ is respectively $(1,1)$ and $(0,0)$; on the matrix block it is
\[
 e=\begin{pmatrix}1&0\\0&0\end{pmatrix},\qquad
 f=\begin{pmatrix}1-z&1\\z(1-z)&z\end{pmatrix}.
\]
In particular $A$ is free of rank ten. At $v=1$ over a field $k$ of characteristic different from 2 and 3,
\[
 A_1\simeq k\times k\times M_2(k[\varepsilon]/\varepsilon^2).
\]
The actual specialization map to the ambient tensor group algebra has kernel the four-dimensional square-zero Jacobson radical and image the diagonal $k[S_3]$.

Over the local characteristic-zero ring, the generic center branches collide to order two:
\begin{equation}\label{eq:r3collision}
 z_1-z_2=\frac{D-1}{D^2}(v-v^{-1})^2.
\end{equation}
The algebra is also the global coherent endomorphism algebra of an explicitly given vector bundle on the ordinary commutative scheme
$S=\Spec R\sqcup\Spec R\sqcup\Spec C$.
\end{finding}
\begin{proof}
The source input is the complete presentation in IV, Theorem~11.16, equations~(111)--(112), printed p.~65. With $x=P_1$, $y=P_2$, $a_1=\delta$ and $a_2=D-2$, it is
\[
 x^2=Dx,\quad y^2=Dy,\quad
 x(yx-a_1^2)(yx-a_2^2)=y(xy-a_1^2)(xy-a_2^2).
\]
Set $e=x/D$, $f=y/D$. These are idempotents, and $z=(e-f)^2=e+f-ef-fe$ is central: multiplying on either side by $e$ gives $e-efe$, and multiplying on either side by $f$ gives $f-fef$. Consequently
\[
 efe=(1-z)e,\qquad fef=(1-z)f.
\]
Expansion of the displayed braid relation gives
\[
 D^5(z-z_1)(z-z_2)(e-f)=0.
\]
For example, its two sides expand using $xyx=D^3(1-z)e$ and $xyxyx=D^5(1-z)^2e$, and the corresponding expressions in $f$; the two factors are $D^2(1-z)-a_i^2$. Thus the presentation is equivalent to the idempotent relations and $g(z)(e-f)=0$, where $g(z)=(z-z_1)(z-z_2)$.

Multiplication by $e-f$ gives $zg(z)=0$. Since $g(0)=z_1z_2$ is a unit, the central Chinese remainder decomposition splits the algebra into its $z=0$ and $g(z)=0$ blocks. In the first, $g(0)(e-f)=0$ forces $e=f$, giving $R[e]/(e^2-e)=R\times R$.

In the second block the central ring is $C$. Every word is a $C$-linear combination of $1,e,f,ef$, by removing repeated idempotents, shortening alternating triples with the preceding identities, and replacing $fe$ by $e+f-ef-z$. The matrices in the statement satisfy those identities. Both $z$ and $1-z$ are units in $C$: their evaluations at the two roots are units, since
\[
 1-z_1=1/D,\qquad 1-z_2=(D-2)^2/D^2.
\]
Equivalently their norms in the quadratic algebra are units, which gives their inverses. Now
\[
 ef(1-e)=E_{12},\qquad (1-e)fe=z(1-z)E_{21}.
\]
Together with $e$ and $1-e$ these generate every matrix unit. The matrices $1,e,f,ef$ are also independent: the lower-left coordinate of a relation first kills its $f$ coefficient; the upper-right then kills its $ef$ coefficient; the lower-right kills its scalar coefficient; the upper-left kills the remaining $e$ coefficient. This proves the isomorphism and freeness.

At $v=1$, $D=4$ and $z_1=z_2=3/4$, giving the dual-number description. Its ideal
$J=0\times0\times M_2(\varepsilon k)$ has square zero and semisimple quotient $k\times k\times M_2(k)$, so it is exactly the Jacobson radical and has dimension four. In the original ambient inclusion,
\[
 P_i=2(1+s_i\otimes s_i).
\]
These generate the diagonal group algebra, since $s_i\otimes s_i=P_i/2-1$. Its dimension is six and it is semisimple in the stated characteristics. The image of $J$ is a nilpotent ideal of that semisimple image and hence is zero; dimension then makes $J$ the entire kernel.

Equation~\eqref{eq:r3collision} follows from $D-4=(v-v^{-1})^2$. Evaluation at the two roots identifies
\[
 C=\{(u_1,u_2)\in R^2:u_1-u_2\in(z_1-z_2)R\}.
\]
Indeed $a+bz$ evaluates to a pair whose difference is $b(z_1-z_2)$; conversely such a pair gives this $b$ and $a=u_1-bz_1$. Injectivity follows because $R$ is a domain and the two generic roots differ. At $v=1$, the factor multiplying $(v-1)^2$ in \eqref{eq:r3collision} is a unit over the local characteristic-zero ring. The generic central idempotents divide by $z_1-z_2$ and do not extend through this collision; the finite algebra $C$ over this local ring has local special fiber and is itself local, so has no nontrivial idempotent.

Finally take the trivial vector bundle $E$ of ranks 1, 1 and 2 on the three components of $S$. Its algebra of global coherent endomorphisms is exactly $R\times R\times M_2(C)$. The displayed matrices retain the named nonstandard generators. This proves the coherent realization.
\end{proof}
This strengthens the construction beyond an abstract matching-dimension model. It also corrects a possible flatness misreading: the family remains rank ten; the \emph{specialized inclusion} loses injectivity. IV, Remark~11.9 (p.~61), already distinguishes abstract specialization from ambient image; Proposition~11.14 (pp.~64--65) gives the generic simple dimensions $1,1,2,2$ and the merging of the two two-dimensional representations.

Over a finite field of good characteristic at $v=1$, $S$ has two reduced rational points and one double rational point. Its geometric \(\ell\)-adic cohomology is that of three rational points, since a nilpotent thickening does not change the \emph{\'etale topos} (\cite[Proposition~59.45.4, tag~03SI]{StacksNil}). Thus it is pure in degree zero and has dimension three; this purity does not see the coherent radical. Replacing the three vector-bundle ranks by constant \(\ell\)-adic sheaves gives a chosen endomorphism algebra of dimension six on the reduced points, not a faithful realization of the ten-dimensional coherent algebra.

A full endomorphism algebra of a finite semisimple object over an algebraically closed coefficient field is a product of matrix algebras; it cannot be $A_1$. An algebra acting as a \emph{subalgebra} of such an endomorphism algebra can be nonsemisimple, so that broader construction is not excluded. The coherent bundle above is not yet the original GCT tensor module, and the scheme model does not identify VIII's basis or formal quantum grading with Frobenius weights. These remain the precise comparison tasks. In Grothendieck's language the result lives in finite-flat schemes and coherent $\End(E)$, not in an unproved RH; EGA~I's affine scheme formalism, SGA~3's module/comodule language, and cote~90, scan~17, supply the framework, while the complete algebra proof is given here.

\subsection{Actual pure classical multiplicity objects via geometric Satake}
\begin{finding}[\fid{F20}: Schur summands have pure multiplicity spaces; their cycle action need not be positive]\label{f:satake}
Let $G=\GL_r/\F_q$ be split, with $\ell\nmid q$, and use the symmetric geometric Satake category with $\overline\Q_\ell$ coefficients. For a dominant partition $\mu$ of length at most $r$, let $\IC_\mu$ be the natural unnormalized Weil intersection complex of its Schubert variety, perverse of weight
\[
 D_\mu=\dim\Gr_\mu=\langle2\rho,\mu\rangle
 =\sum_{i=1}^r(r+1-2i)\mu_i.
\]
For $\nu\vdash d$, a primitive Young idempotent $e_\nu\in\Q[S_d]$ acts, through the parity-corrected symmetric commutativity, on
$K=\IC_\mu^{\star d}$. Its image $P_{\nu,\mu}$ is a pure perverse Weil sheaf of weight $dD_\mu$. The canonical multiplicity Weil module
\[
 M_{\lambda;\nu,\mu}
 =\Hom_{\mathrm{Sat}_{\overline{\F}_q}}
 (\IC_\lambda,P_{\nu,\mu})
\]
is pure of weight $dD_\mu-D_\lambda$ and has dimension
\[
 \dim M_{\lambda;\nu,\mu}
 =[S_\lambda(\overline\Q_\ell^r):
   S_\nu(S_\mu(\overline\Q_\ell^r))].
\]
On the product Grassmannian for $\GL_m\times\GL_n$, apply the Schur idempotent for $\nu\vdash d$ to the $d$th convolution power of
$\IC_{(1),m}\boxtimes\IC_{(1),n}$. Its multiplicity module at
$\IC_{\lambda,m}\boxtimes\IC_{\mu,n}$ has dimension $g(\lambda,\mu,\nu)$ and is pure of weight
\[
 d(m+n-2)-D_{\lambda,m}-D_{\mu,n}.
\]
These are ordinary geometric realizations of the classical coefficients, including rectangular cases within the dimension bounds. They do not by themselves identify the nonstandard quantum bases.

Even in this construction the symmetric-group action on a geometric multiplicity space need not permute a basis. For $G=\GL_2$, $\mu=(1)$ and $d=2$, the determinant multiplicity line is $\overline\Q_\ell(-1)$ and the transposition acts by $-1$. Thus no basis of that line is an $S_2$-permutation basis.
\end{finding}
\begin{proof}
The literature inputs are geometric Satake, its finite-field convolution/symmetry, and ordinary weight theory. We use the published version of Richarz \cite[Theorem~3.1, p.~214; proof pp.~226--227; Proposition~4.1, p.~227; Lemma~5.4 and its proof, p.~232]{Richarz}. The tensor equivalence is also Mirkovi\'c--Vilonen, Theorem~7.3 \cite{MV}. These theorems are cited, not re-proved here.

The twisted product of $d$ Schubert varieties is projective and is smooth-locally the product for purposes of its intersection complex. Thus its IC is the twisted external product of the $\IC_\mu$ and is pure of weight $dD_\mu$; this is the IC identification in the proof of Richarz's Lemma~5.4, together with BBD~5.3.2. Its proper convolution pushforward is pure of that same complex weight (BBD, Remark~5.4.9) and is perverse (Richarz, Theorem~3.1). This proves the assertion about $K$.

The symmetry constructed over $\F_q$, with the component-parity sign specified in the cited proof, gives an $S_d$ action by Weil endomorphisms. In the abelian perverse category the idempotent splits; its image is a direct summand, hence remains pure by BBD~5.3.1. Under the symmetric monoidal Satake equivalence, $\IC_\mu$ corresponds to $S_\mu(\overline\Q_\ell^r)$ and $K$ to its $d$fold tensor power. In Schur--Weyl decomposition, the primitive idempotent $e_\nu$ has rank one on the Specht factor $[\nu]$ and vanishes on the other factors. Its image therefore corresponds to one copy of the Schur functor $S_\nu(S_\mu(\overline\Q_\ell^r))$. A central isotypic idempotent would instead give $\dim[\nu]$ copies; that is why the primitive condition matters.

Here is the Frobenius argument for the multiplicity weight. Geometric semisimplicity gives the canonical evaluation isomorphism
\[
 \bigoplus_\lambda\IC_\lambda\otimes
 \Hom(\IC_\lambda,\overline P_{\nu,\mu})
 \longrightarrow\overline P_{\nu,\mu}.
\]
It is the canonical isotypic decomposition, not an arbitrarily chosen decomposition. Every support and natural $\IC_\lambda$ is defined over the split base $\F_q$, so Frobenius preserves each summand. Give its Hom factor the conjugation Frobenius induced by the two fixed Weil structures. Evaluation is then Weil-equivariant, and the isotypic summand is a pure perverse Weil sheaf of weight $dD_\mu$. This descent/subquotient formalism is BBD~5.1.2; purity of the summand is BBD~5.3.1. On its own smooth open Schubert stratum, $\IC_\lambda$ is the constant sheaf shifted by $D_\lambda$. Restriction to this open preserves both weight inequalities because $j^*=j^!$ (BBD~5.1.14). The Hom module shifted by $D_\lambda$ is therefore pure of weight $dD_\mu$, which means the module itself has weight $dD_\mu-D_\lambda$ by the shifted convention of BBD~5.1.8. This proves the all-eigenvalue purity statement without assuming arithmetic semisimplicity. The tensor equivalence identifies its dimension with the required multiplicity.

For the product group, external product gives the simple object corresponding to $V\boxtimes W$. Its weight is $(m-1)+(n-1)$. The same Schur-idempotent construction corresponds to $S_\nu(V\otimes W)$ as a $\GL_m\times\GL_n$ representation. The definition of the Kronecker coefficient and the just-proved Hom-weight calculation give the stated formula.

Finally consider $\GL_2$ and two minuscule steps. The twisted convolution variety is a smooth surface. Over the central lattice $t\mathcal O^2$, its fiber consists of the intermediate lattices between $t\mathcal O^2$ and $\mathcal O^2$, hence of lines in $\mathcal O^2/t\mathcal O^2$; it is $\PP^1$. The contribution to the perverse degree-zero stalk is $H^2(\PP^1)=\overline\Q_\ell(-1)$. To identify it as the determinant multiplicity, Satake gives
\[
 V\otimes V=\Sym^2V\oplus\wedge^2V,
\]
with determinant multiplicity one. Equivalently the non-open relevant fiber supplies this unique one-dimensional summand; the IC of the open two-dimensional stratum has no perverse degree-zero stalk contribution at the boundary point by the intermediate-extension condition. The corrected symmetric commutativity corresponds to the ordinary tensor swap, which sends
$e_1\otimes e_2-e_2\otimes e_1$ to its negative. The sole nonzero basis vector of a one-dimensional permutation representation would have to be fixed, so no such basis exists. The symmetric projector kills this line, while the alternating projector keeps it.
\end{proof}
This is a constructive reduction of the ordinary-geometry gate for \emph{classical} coefficients. Its purity uses the existing Weil/BBD theorems; it is not a new application of \fid{F03}, and it does not provide the path-hypersurface identities for a nonstandard canonical coefficient. The remaining problem is more specific: describe the actual projector image, its graded basis and its compatibility with the designated quantum operations, together with effective labels when a counting rule is required. The $\GL_2$ example prevents replacing that task by unsigned orbit-counting on the unprojected component basis.

The construction also provides an exact Grothendieck-framework dictionary. Convolution is a proper constructible pushforward; the Young idempotent is an actual projector in a characteristic-zero symmetric tensor category; its image is an effective object whose $K_0$ class is the Schur operation. The weights record the unnormalized IC shifts, so no half-Tate convention is hidden. The motivating language of projectors occurs in cote~16, scans~35, 37 and~39; the actual theorem here uses the modern Satake equivalence and the precisely cited SGA/BBD weight formalism, not an attribution to those leaves. The explicit symmetric-group action is the additional datum beyond virtual $\lambda$ identities.

\section{What is reached in the GCT cases requested}\label{sec:rectangles}
We keep three different notions separate: scalar nonnegativity of a representation multiplicity; a positive effective rule for it; and positivity of normalized canonical structure constants. The first is already part of semisimple representation theory and \fid{F12}. The others do not follow from it.

Write $n\times d=(d^n)$, $\lambda\vdash nd$, and
\[
 g_n(\lambda)=g(\lambda,(d^n),(d^n)),\qquad
 a_\lambda(d[n])=[S_\lambda V:\Sym^d(\Sym^nV)].
\]
Let $q^m_\lambda(d[n])$ denote the degree-$d$ multiplicity in the coordinate ring of the orbit closure of $X_{11}^{n-m}\operatorname{per}_m$, in IP's covariant convention. The known inequality is $0\le q^m_\lambda(d[n])\le a_\lambda(d[n])$ (IP, equation~(1.3)).

\subsection{Exact rectangular results already proved}
The attached Ikenmeyer--Panova paper \cite{IP}, matching arXiv 1512.03798v2, states:
\begin{itemize}
\item Theorem~1.4: $q^m_\lambda(d[n])>g_n(\lambda)$ implies $\operatorname{dc}(\operatorname{per}_m)>n$.
\item Theorem~1.6: if $n>3m^4$ and $g_n(\lambda)=0$, then $q^m_\lambda(d[n])=0$.
\item Theorem~1.8: $q^m_\lambda(d[n])>0$ requires $\ell(\lambda)\le m^2$ and $|\bar\lambda|\le md$, where $\bar\lambda$ removes the first row.
\item Theorem~1.10: put
$\mathcal X=\{(1),(1,1),(1^4),(1^6),(2,1),(3,1)\}$.
If $\bar\lambda\in\mathcal X$, then $a_\lambda(d[n])=0$. If $\bar\lambda\notin\mathcal X$ and there is $m\ge3$ with
\[
 \ell(\lambda)\le m^2,\quad |\bar\lambda|\le md,\quad
 d>3m^3,\quad n>3m^4,
\]
then $g_n(\lambda)>0$.
\end{itemize}
These are substantive existing positivity and no-occurrence results. They are not new consequences of our conditional sign theorem. Revised IV, Theorems~1.6/15.21(vii)--(viii), also supplies a positive highest-weight crystal rule when the two restricted factors have at most two rows (and the $\GL_4$ label at most four parts). That includes two-row rectangles, not arbitrary-height rectangles in the asymptotic regime.

\subsection{Occurrence versus multiplicity}
Let $\Omega_n=\overline{\GL_{n^2}\cdot\det_n}$, $Z_{n,m}$ the padded-permanent orbit closure, and $K_n(\lambda)$ the degree-$d$ multiplicity in $\C[\Omega_n]$. B\"urgisser--Ikenmeyer--Panova \cite[Theorem 1.4]{BIP}, in the attached arXiv 1604.06431v3, proves for $n\ge m^{25}$:
\[
 q^m_\lambda(d[n])>0\quad\Longrightarrow\quad K_n(\lambda)>0.
\]
Theorem~1.5 is the numbering in an older version, not this attachment. The paper's \S1(b) distinguishes
\[
 K_n(\lambda)\le sk_n(\lambda)\le g_n(\lambda),
\]
where $sk_n$ is the symmetric rectangular coefficient controlling the determinant \emph{orbit}, not automatically its closure.

If $Z_{n,m}\subseteq\Omega_n$, restriction is a graded equivariant surjection
$\C[\Omega_n]\twoheadrightarrow\C[Z_{n,m}]$. Semisimplicity of the characteristic-zero representation in each finite-dimensional degree implies
\[
 q^m_\lambda(d[n])\le K_n(\lambda).
\]
Thus a strict reverse multiplicity inequality remains a valid obstruction, as BIP \S1(c) explains. A theorem about nonzero support does not settle it. The ordinary Kronecker upper bound supplies a sufficient stronger inequality, not an equality with closure multiplicities.

\paragraph{Consequence of the present work.}
\fid{F03}--\fid{F04} establish purity only for data satisfying their trace and sign hypotheses. No such data have been identified here for the normalized coefficients in VIII, the full rectangular region, or the closure multiplicity comparison. Therefore no new general positivity or multiplicity obstruction is derived from them. This is the exact limit of the proved transfer, not a restatement of numerical nonnegativity as a new theorem. The permanent--determinant problem concerns algebraic complexity classes; an additional argument would be required before making a Boolean $P\ne NP$ claim.

\part{The findings in Grothendieck's language}
\section{An exact dictionary, not an attribution of new theorems}\label{sec:groth}
The new deductions above are restated here in the requested frameworks. A framework locator says where the language and operations live; it does not claim that a theorem about GCT occurs in Grothendieck's manuscripts. Where no archival proof or identification was verified, we say so explicitly.

\subsection{Schemes and group schemes over the integers}
The classical objects are group schemes $\GL_{r,\Z}$ and their polynomial representations on finite locally free modules. Tensor products and integral Schur constructions give maps such as
\[
 \GL(V)\times\GL(W)\longrightarrow\GL(V\otimes W),\qquad
 \GL(V)\longrightarrow\GL(S_\mu V).
\]
Representations of an affine group scheme are comodules of its \emph{commutative} coordinate Hopf algebra (SGA~3, I, \S4.2 and 4.7.2). Representations of diagonalizable groups correspond to graded modules (I, 4.7.3). These descriptions carry \eqref{eq:plethysm}, \fid{F11} and \fid{F13} into group-scheme language. Cote~90, scan~17, gives an authentic manuscript anchor for a flat affine group over a base, a group action, and equivariant quasi-coherent modules.

Over characteristic zero the representation categories of the reductive groups here are semisimple, so their actual classes form the positive monoid in the irreducible basis of $K_0$. This is the precise content of scalar nonnegativity. Over $\Z$ or in characteristic $p$, do not replace a filtration by a direct sum of simple modules. A divided Young symmetrizer is only an idempotent when its denominators are invertible; integral Schur functors must be constructed integrally. For the integral Specht-model representations used in the Kronecker and plethysm instances of \fid{F12}, tensor the free integral lattices with $\Z[1/|H|]$. The averaging idempotent then has a finite locally free direct-summand image, an effective $K_0$ class whose complex rank is the stated multiplicity. More generally this construction works over any coefficient ring carrying the representation and in which $|H|$ is invertible; an arbitrary complex representation need not descend to $\Z[1/|H|]$.

The nonstandard bialgebra preserves some comodule language but is not automatically a commutative affine group scheme. \fid{F11}'s matrix calculation proves why a reduction modulo $p$ supplies no automatic absolute Frobenius. This is a limitation of that construction, not a statement against separately defined quantum Frobenius maps.

\subsection{Representation rings, $\lambda$-rings and the $\gamma$-filtration}
For locally free sheaves, $\lambda^i[E]=[\wedge^iE]$ and Schur operations are the associated symmetric-function operations. On virtual classes the identity $\lambda_t(x-y)=\lambda_t(x)/\lambda_t(y)$ belongs to the group-completed ring. SGA~6, Expos\'e~V, \S\S2--3, contains the $\lambda$ and $\gamma$ formalism; cote~133, scan~18, explicitly displays symmetric/exterior generating functions. Cote~97, scan~4, concerns Chern classes of representations. These are the appropriate anchors for the formal calculation in \fid{F11}.

The representation ring example is an exact dictionary, not an analogy: $z$ is the standard one-dimensional comodule of $\mathbb G_m$; $z-1$ is a rank-zero virtual class; its exterior-power series gives $\lambda^2(z-1)=1-z$. The augmentation ideal and its $\gamma$-filtration retain characteristic-class information but introduce no order turning this virtual class into an actual representation. In particular the equality of two expressions in a $\lambda$-ring cannot replace the effective-simple-basis hypothesis of \fid{F07}.

For singular schemes, one must distinguish $K_0$ of vector bundles or perfect complexes from $G_0$ of coherent sheaves. The GCT graded pieces themselves are finite-dimensional representations, so their representation-$K_0$ classes are well-defined without identifying $K_0(X)$ and $G_0(X)$ on a singular orbit closure. Coherent Euler pushforward may involve alternating higher cohomology; an actual section module is more data than a virtual Euler class.

\subsection{Weil realizations and virtual versus effective traces}
In the trace formalism, a bounded complex of finite-dimensional Weil representations yields
\[
 Z(K,T)=\prod_i\det(1-TF\mid H^i(K))^{(-1)^{i+1}},\qquad
 \log Z(K,T)=\sum_{n\ge1}\Tr(F^n\mid K)\frac{T^n}{n}.
\]
The trace is an alternating trace. \fid{F02} bounds the surviving finite spectrum after equal roots have been grouped. \fid{F05}(b) says that $K_0$ and $Z(K,T)$ can erase identical spectra in adjacent degrees. An effective parity statement prevents this erasure; that is the exact role of (S2).

The geometric dictionary of \S\ref{sec:tracegeom} is the proof-carrying transfer of \fid{F01}--\fid{F04} and \fid{F17}. Character orthogonality is expressed by Artin--Schreier sheaves; the external tensor product realizes powers; compactly supported direct image realizes the sum; one-degree support fixes the alternating sign; Tate twist supplies the factor $q^n$; duality supplies the reciprocal eigenvalues. Framework locators are SGA~4, XVII, \S3 (compact supports), XVIII, 3.2.5 (smooth duality); SGA~5, III, 4.7 (Lefschetz--Verdier); and SGA~$4\frac12$, \emph{Rapport}, 3.2 (the trace formula). The exact ordinary weight estimates are Weil~II, 3.3.1 and 3.3.4--3.3.6. These general theorems are cited background; this run does not re-prove their machinery.

For a scheme over $\F_q$ and a closed point $x$ with residue field $\F_{q^{\deg x}}$, a Weil sheaf is pure of weight $w$ if every eigenvalue $\alpha$ at $x$ has
\[
 |\sigma(\alpha)|=q^{(\deg x)w/2}
\]
for every relevant complex embedding. For complexes the weight conditions are degree-shifted and involve duality; purity of a complex does not assert that every singular stalk has exact weight equal to its degree. The KL theorem uses the stronger stalkwise even/Tate statement \eqref{eq:KL}. This distinction translates \fid{F06} and prevents a false purity-to-coefficient inference.

\subsection{Orbit closures, Hilbert/Quot, and motives}
\fid{F15}'s kernel construction is the affine scheme-theoretic-image construction of EGA~I (1960), \S9.5. Its explicit $f=p$ example is already a group-scheme example over $\Spec\Z$. Projective closures with fixed Hilbert polynomial and coherent quotients belong to FGA, Bourbaki Expos\'e~221, \S3, Theorems~3.1--3.2, pp.~259--260. Fixing the Hilbert polynomial uses the chosen polarization, and a Quot functor specifies the coherent sheaf being quotiented. They do not turn coherent section multiplicities into intersection-cohomology multiplicities by representability alone.

For smooth projective schemes with algebraic correspondences, pure motives and their realizations provide another language for duality and projectors. The Tate realization in \fid{F16} is legitimate. By contrast a general singular orbit-closure intersection complex is not automatically a Chow motive furnished by SGA~4/5/6, and a quantum coordinate algebra acquires no such motive from its name. Cote~16, scans~35, 37 and~39, supplies a historical anchor for Lefschetz operators and primitive projectors; scans~47--49 discuss conditional correspondence/motive consequences, with conjectural letters interpreted in the letter's own notation. Their algebraic-projector and index assumptions cannot be replaced by the Euler sign $(-1)^i$.

\subsection{Finding-by-finding transfer table}
{\small\begin{longtable}{p{.12\linewidth}p{.48\linewidth}p{.32\linewidth}}
\toprule IDs&Exact Grothendieck-framework form&Locator/anchor and limit\\\midrule\endhead
F01--F04, F17&Trace of external tensor powers on the path hypersurface; parity of $R\Gamma_c$; Tate duality; optional finite base extension for each $k$. The calculations of \S\ref{sec:main} apply after the trace-formula identification.&SGA~4 XVII \S3, XVIII 3.2.5; SGA~5 III 4.7; Weil~II 3.3.1. Cote~16 provides only the distinct projector/duality framework, not this theorem.\\
F05&The zero virtual class $[\Q_\ell(-1)]-[\Q_\ell(-1)]$ has invisible large eigenvalues. A pure arithmetic Weil module can be a nontrivial extension.&Triangulated $K_0$; BBD 5.3.8 and its following example, pp.~138--139. No archival proof of the new examples is asserted.\\
F06--F07&Effective graded simple classes yield nonnegative multiplicities; a positive phase and Lefschetz maps are additional. Euler classes alone do not suffice.&BBD 5.3.2, 5.4.5, 5.4.10; cote~16, scans~35, 37, 39 for operator language.\\
F08&The rank vector bundles on two points form a special $\lambda$-ring; an effective integral basis need not be its simple positive basis.&SGA~6 V \S\S2--3; cote~133, scan~18. The entire counterexample is proved here.\\
F09--F10&An algebra inclusion need not preserve the effective monoid of a split $K_0$ categorification; exclusive monomial pivots give a sufficient alternative.&Split Grothendieck groups, effective objects, grading shifts; cote~90, scan~17 for equivariant-module context. We do not attribute the nonstandard inclusion claim to the manuscripts.\\
F11&Virtual exterior powers fail effectivity; noncommutative $p$th powers fail additivity, so do not define a scheme Frobenius.&SGA~3 I 4.2, 4.7.2--3; SGA~6 V \S\S2--3; cotes~97/4 and~133/18.\\
F12--F13&For the integral Specht models, invariant projectors are direct summands over $\Z[1/|H|]$; positive restriction matrices compose in representation $K_0$.&SGA~3 I 4.7.2--3; cote~90/17. Actual theta/GCT functor is still absent.\\
F14&Frobenius-stable subquotients inherit a weight; endomorphism spectra are ratios.&Weil representations; BBD 5.3.1 for perverse subquotient weight bounds. Realization/stability is still required.\\
F15&Kernel-defined orbit schemes over $\Z$; coherent polarization and linearization retained in $K_0/G_0$; finite-field base change is separate.&EGA~I \S9.5; FGA~221, 3.1--3.2; cote~90/17 for group actions.\\
F16&The centered polynomial of $\mathbf1\oplus\mathbb L^{\oplus2}\oplus\mathbb L^2$ realizes the rank-one product coefficient; its action comparison fails.&Projective bundle/K\"unneth calculations; cote~16 correspondence language. Pure Tate pieces, not an ungraded pure weight-zero module.\\\bottomrule
\end{longtable}}

\section{Archival audit and the Fable leads}\label{sec:archive}
The original Montpellier PDFs add a cover. For the documents inspected here, a bottom-stamped archive scan $s$ is one-based PDF page $s+1$. Catalogue titles and scanned contents are distinguished. The following checks were visual readings of the specified leaves, not certifications of full manuscript proofs.

{\small\begin{longtable}{p{.19\linewidth}p{.36\linewidth}p{.37\linewidth}}
\toprule Cote and scan&Verified content&Consequence for this investigation\\\midrule\endhead
90, scan 17 (PDF 18)&Flat affine group over a base; group action; quasi-coherent equivariant modules.&Authentic anchor for the representation/group-scheme dictionary.\\
97, scan 4 (PDF 5)&Chern classes of regular representations and products over characters.&Characteristic-class anchor, not a plethysm-positivity theorem.\\
133, scan 18 (PDF 19)&Symmetric and exterior generating operations, including $\sigma_t(E)\lambda_{-t}(E)=1$.&Direct language anchor for the $\lambda$-ring calculation.\\
133, scans 19--21 (PDF 20--22)&Virtual finite $G$-sets and operations; scan 20 explicitly names plethysm.&These finite-set operations must not be identified with a proved representation-theoretic positive rule.\\
16, scans 35, 37, 39 (PDF 36, 38, 40)&Handwritten Lefschetz/primitive-summand/projector formulas, including hyperplane-section comparisons.&Relevant standard-conjecture framework. No identification with Alp\"oge's path sign is proved.\\
16, scans 34, 36, 38, 40&Reused typescript backs; scan 34 starts Karoubi's \emph{Cohomologie des cat\'egories de Banach}.&The Fable range ``34--40'' must be refined; these backs are not Grothendieck's handwritten Hodge-index proof.\\
16, scans 47--49 (PDF 48--50)&Letter to Coates: correspondences/duality; on scan 48, $C(X)$ plus $I(X\times X)$ implies correspondence-ring semisimplicity; on scan 49, deduction of Weil from $C$ and $A$, in the letter's notation.&Conditional standard-conjecture statements. The semisimplicity passage is scan 48, not scan 50.\\
16, scan 50 (PDF 51)&Separate handwritten endomorphism-algebra lemma.&Not the location of the preceding semisimplicity sentence.\\
1, proposed scans 79--80, 111--112&Not successfully verified in this run.&The proposed $\Xi$/operator-determinant identification is not used as a result.\\\bottomrule
\end{longtable}}
Cote~133's catalogue title concerns Witt groups, quadratic forms and formal groups; ``plethysm'' is a content lead, not its official title. Cote~60 was found in the catalogue and inspected at thumbnail level, but no precise theorem from it is used.

The Fable suggestions are treated as follows. Suggestion~3's norm estimate is proved and extended in \fid{F01}; its literal positive-Fourier-operator identification fails by \fid{F05}(e). Suggestion~4's identification of the cohomological sign with a polarized evanescent intersection form remains unproved: parity in a trace formula and a positive-definite form are different structures. Suggestion~8 is made precise by \fid{F13}, conditional on an actual functor and endpoint identification. No assertion about the unprovided earlier runs or their theorem IDs is imported.

Some published-source checks have a narrower scope than the archive checks. The original SGA~6 scan could not be visually retrieved here; Expos\'e~V, \S\S2--3, is carried as background. The specific locator V, Lemma~2.2.1, p.~315, is independently cited by Altman--Kleiman \cite[\S2.10]{AK}; it is a locator check, not a re-proof of SGA~6. The modern transcription consulted as a finding aid is not used as technical authority. All elementary $\lambda$ identities used in the paper have their displayed direct proofs.

\part{Unproved bridges, attempted routes, and final state}
\section{What a sufficient nonstandard package would have to contain}\label{sec:gates}
There are no unchecked affirmative ``new theorems'' in this part. The following is a sufficient research specification, not a quotation from Mulmuley and not a proved construction.

\begin{enumerate}
\item \gate{R: realization and normalization.} For each relevant nonstandard coefficient construct a finite-field Weil object or a defined categorical replacement, with explicit Frobenius, trace formula, degrees and weights. Identify $v$ with the correct grading, and specify signs, half-powers and rational rescalings.
\item \gate{S: spectral sign and duality.} Verify the pointwise input bound, the finite-spectrum identification, opposite parity for unbounded $k$ (over finite extensions if necessary), and controlled background terms. Supply either the Tate self-duality of \fid{F03} or the dual system of \fid{F17}. These inputs imply the stated purity theorem by the checked proof.
\item \gate{E: effective basis and operations.} Identify the normalized coefficient with a graded multiplicity of actual simple objects and prove compatibility with multiplication/comultiplication. Purity alone does not give this. The fixed inherited ambient cone is ruled out by \fid{F09}; the new rank-one cone of \fid{F16} changes the action.
\item \gate{L: Lefschetz and defects.} Prove the appropriate Lefschetz maps for unimodality. In general plethysm, define the nonflatness correction and prove the quantitative error bounds of the intended version of VIII 2.9/3.6. A spectral error bound such as \fid{F04} is not automatically that coefficient-vector error bound.
\item \gate{C: canonical and comparison data.} Prove basis existence/globalization, cell and quasi-cell claims, and the required specialization. For general plethysm retain VII's duality/reciprocity hypotheses; for Kronecker do not re-open the already proved revised-IV duality.
\item \gate{P: effective labels and polyhedra.} Produce polynomial-length witnesses and a polynomial-time checking predicate for a \(\#P\) interpretation. PH1 additionally requires its Ehrhart family, dimension/bitlength bounds and separation oracle; SH or the relevant index bound is separate.
\item \gate{O: actual obstructions.} Establish a strict multiplicity comparison for the orbit closures in an appropriate asymptotic family. Positive formulas and efficient decisions do not themselves produce that inequality.
\end{enumerate}
The theorem proved here removes a \emph{spectral estimate} once R and the explicit S hypotheses are supplied. It does not construct R, E, L, C, P or O. Replacing all seven by the single phrase ``nonstandard RH implies positivity'' hides independent claims.

For completeness, an effective-label claim has a simple formal implication: if the labels $\alpha$ in \eqref{eq:comparison} and witnesses for $m_\lambda^\alpha,n_\pi^\alpha$ have polynomial-length unique encodings and polynomial-time tests, then triples consisting of a label and the two witnesses count the sum of products. A nondeterministic machine guesses and checks the triple, giving a \(\#P\) function. This is just the definition of \(\#P\); neither the encodings nor PH1 follows from the abstract multiplicity equality. No complexity-class collapse or lower bound is asserted.

\section{Approaches attempted and their exact outcomes}
{\small\begin{longtable}{p{.23\linewidth}p{.40\linewidth}p{.29\linewidth}}
\toprule Approach&Proved outcome or failure point&Next unresolved input\\\midrule\endhead
Literal source formulation&No finite-field nonstandard RH statement in the supplied versions; actual signed/near-positive targets located.&Choose a precise R/E target.\\
Path orthogonality and repeated Parseval&F01 gives the bound, but the constant trace example defeats purity without spectral signs.&S1/S2 realization.\\
Tensor powers and inverse roots&F02--F04 remove the fixed exponent loss; F17 allows separate finite extensions and a dual system.&Prove the concentration/error and duality hypotheses.\\
Fourier-positive operator&The bare Fourier-after-nonnegative-multiplication claim is false, F05(e).&A different operator and a proved identity to the coefficient.\\
Ordinary ambient purity&Stable subquotients inherit purity, F14. The actual inherited positive Hecke cone has only scalar image, F09.&A new effective realization/heart or different character map.\\
Change of rank-one positive model&F16 gives positive Tate structure data for the abstract algebra.&Restore the original tensor action; the sign-tensor-sign test shows it currently fails.\\
Positive supports/crystals&Monomial pivots suffice, F10. Global positive disjoint supports and local disjointness alone fail.&Global exclusive pivots or another valid basis argument.\\
Orbit/Hilbert/Quot geometry&Integral models and ordinary IC purity exist; polarized coherent multiplicities are not determined by underlying Weil cohomology, F15.&A coherent-to-constructible multiplicity comparison.\\
Determinants and Gram forms&Positive projectors already express the ranks, F12; chosen-basis coefficient positivity fails in F08.&Image basis, independence and operation compatibility.\\
Theta and interlacing chains&F13 gives path/polytope counts for specified exact functors.&Identify the actual theta chain and endpoints with GCT restriction.\\
$\lambda$-rings and motives&Formal operations and Tate examples translate exactly; virtual effectivity and algebraic projectors remain extra.&Effective classes and the required realization, not only ring identities.\\\bottomrule
\end{longtable}}

\section{States of the five requested tasks}
\begin{enumerate}
\item \textbf{Statement needed.} The source-location and target audit is completed as A01 and \S\ref{sec:targets}. The demand for exact nonstandard Frobenius objects/weights cannot be met from those papers because they are not specified. The replacement package in \S\ref{sec:gates} is explicitly our sufficient specification.
\item \textbf{Template.} The standard geometric dictionary and every missing role are stated in \S\ref{sec:KL}. An actual nonstandard replacement remains open at gate R/E.
\item \textbf{Sign-only mechanism.} F01--F04 and F17 are proved and checked; abbreviated versions are refuted by F05. Application to the designated nonstandard objects remains open at R/S. The specific inherited-positive-category route is refuted by F09, while F16 shows the scope of that refutation.
\item \textbf{Positivity reached.} No new general canonical-basis, Kronecker/plethysm counting, or closure-multiplicity theorem follows from the available data. The exact known rectangular and two-row results are retained with their locators. General requested positivity remains open at E/L/C/P; multiplicity obstruction construction remains at O.
\item \textbf{Theta positivity.} F12--F13 state and prove the exact available projector/path consequences. The theta/GCT functor and basis identification remain open. F08 refutes the generic determinant-to-basis-positivity shortcut.
\end{enumerate}

The final package is therefore marked \textbf{INCOMPLETE}. This is not a claim that no future approach can work. It preserves all checked deductions, the failed shortcuts and the remaining gates. No result here proves or refutes the intended full nonstandard positivity program, or $P\ne NP$.

\appendix
\section{Verification and reproducibility}
The independent spectral audit reconstructed F01, reproved F02 by Ces\`aro orthogonality rather than poles, and checked the two-limit argument of F03. The initial all-conjugate duality gap was recorded and repaired before the final theorem. Later full-paper reviews checked each proof, including the actual Hecke-cone example, the positive-pivot countermodels, the categorical and theta statements, and the scheme/base-change examples. Exact dates, reviewer runs, model identity, exposed effort information and check scopes are in \texttt{CHECK\_EVENTS.tsv}.

The file \texttt{check\_path\_bound.py} checks 54 small scalar/vector examples, including complex weights, for the exact path identity and the estimate. These floating-point checks are supplementary; the proof is F01. The optional exact Murnaghan--Nakayama script and character certificate compute a small rectangular coefficient and check character orthogonality; they are not used in any proof here and are not a Lean certificate. The source manifest records the unchanged attachments by SHA-256. Source reads and whole-proof checks are distinguished from cited background theorems and from formal verification.

The dated ledger records actual changes of gates. Scheduled checkpoints are due only at 45-minute intervals; scheduled progress reviews begin only after approximately three hours. The metadata records which such deadlines, if any, occurred; no review or checkpoint is invented to imply elapsed work that did not occur.

\begin{thebibliography}{99}
\bibitem{Alpoge} L.~Alp\"oge, \emph{Remark}, supplied four-page note, 20 November 2026; \url{https://alpo.ge/remark.pdf}. Main locators: p.~3 footnote~3 and steps~3--5; p.~4 steps~7--8. Supplied formalization entry points: \texttt{ALPOGE\_Main.lean}, \texttt{ALPOGE\_Statements.lean}; scan record \texttt{ALPOGE\_LEAN\_SCAN.tsv}. Results accepted as stipulated; missing imported development not executed.
\bibitem{IV} J.~Blasiak, K.~D.~Mulmuley and M.~Sohoni, \emph{Geometric Complexity Theory IV: nonstandard quantum group for the Kronecker problem}, arXiv:cs/0703110v4, 2013. \url{https://arxiv.org/abs/cs/0703110v4}. Exact statement locators are given in the text; attached source \texttt{Apr13KroneckerGCT4.tex}.
\bibitem{VI} K.~D.~Mulmuley, \emph{Geometric Complexity Theory VI: the flip via saturated and positive integer programming in representation theory and algebraic geometry}, arXiv:0704.0229v4, 2009. \url{https://arxiv.org/abs/0704.0229v4}.
\bibitem{VII} K.~D.~Mulmuley, \emph{Geometric Complexity Theory VII: Nonstandard quantum group for the plethysm problem}, arXiv:0709.0749v2, 2008. \url{https://arxiv.org/abs/0709.0749v2}.
\bibitem{VIII} K.~D.~Mulmuley, \emph{Geometric Complexity Theory VIII: On canonical bases for the nonstandard quantum groups}, arXiv:0709.0751v2, 2008. \url{https://arxiv.org/abs/0709.0751v2}.
\bibitem{IP} C.~Ikenmeyer and G.~Panova, \emph{Rectangular Kronecker coefficients and plethysms in geometric complexity theory}, arXiv:1512.03798v2, 2017; especially Theorems~1.4, 1.6, 1.8, 1.10. \url{https://arxiv.org/abs/1512.03798v2}.
\bibitem{BIP} P.~B\"urgisser, C.~Ikenmeyer and G.~Panova, \emph{No occurrence obstructions in geometric complexity theory}, arXiv:1604.06431v3, 2018; Theorem~1.4, \S1(c), Proposition~3.3, equation~(4.1), Proposition~4.5, Theorem~4.7. \url{https://arxiv.org/abs/1604.06431v3}.
\bibitem{KL} D.~Kazhdan and G.~Lusztig, \emph{Schubert varieties and Poincar\'e duality}, in \emph{Geometry of the Laplace Operator}, Proc. Sympos. Pure Math.~36, 1980, pp.~185--203; Theorem~4.3. The exact normalized form used here is checked through Haines's primary proof below.
\bibitem{Haines} T.~J.~Haines, \emph{A proof of the Kazhdan--Lusztig purity theorem via the decomposition theorem of BBD}; Theorems~2.0.1, 5.0.4, Proposition~3.0.2, \S4. \url{https://www.math.umd.edu/~tjh/KL_purity1.pdf}.
\bibitem{WeilII} P.~Deligne, \emph{La conjecture de Weil II}, Publ. Math. IH\'ES~52 (1980), pp.~137--252; Definitions~1.2.2 and~6.2.2--6.2.4, Theorem~3.3.1 and Corollaries~3.3.4--3.3.6. \url{https://www.numdam.org/item/PMIHES_1980__52__137_0/}.
\bibitem{BBD} A.~Beilinson, J.~Bernstein and P.~Deligne, \emph{Faisceaux pervers}, Ast\'erisque~100 (1982). Exact locators: 5.3.1--5.3.2, 5.3.8 and the following example, 5.4.4--5.4.5, Remark~5.4.9 (p.~143), 5.4.10 (p.~144), 6.2.5. \url{https://www.numdam.org/item/AST_1982__100__1_0/}.
\bibitem{Molev} A.~I.~Molev, \emph{Gelfand--Tsetlin bases for classical Lie algebras}, Theorem~2.1, equation~(2.2), patterns~(2.4); author-hosted preprint, printed pp.~7--9. \url{https://www.maths.usyd.edu.au/u/pubs/publist/preprints/2002/molev-18.pdf}.
\bibitem{Milne} J.~S.~Milne, \emph{Lectures on \'Etale Cohomology}, v2.21, 2013; Example~16.3, p.~109; Theorem~23.2, pp.~139--140. \url{https://www.jmilne.org/math/CourseNotes/LEC.pdf}.
\bibitem{EGA} A.~Grothendieck and J.~Dieudonn\'e, \emph{\'El\'ements de g\'eom\'etrie alg\'ebrique I: Le langage des sch\'emas}, Publ. Math. IH\'ES~4 (1960), \S9.5. Framework locator, with the affine instance proved in F15. \url{https://www.numdam.org/item/PMIHES_1960__4__5_0/}.
\bibitem{SGA3} M.~Demazure and A.~Grothendieck (eds.), \emph{Sch\'emas en groupes (SGA 3)}, Expos\'e~I by M.~Demazure, \S4.2, 4.7.2--4.7.3. Checked in the Gille--Polo re-edition, respecting explicitly marked editorial additions. \url{https://wstein.org/sga/SGA3/Exp1-20jan2010.pdf}.
\bibitem{SGA4} M.~Artin, A.~Grothendieck and J.-L.~Verdier (eds.), \emph{Th\'eorie des topos et cohomologie \'etale des sch\'emas (SGA 4)}, Expos\'es~XVII, \S3, and XVIII, 3.2.5. Background framework references.
\bibitem{SGA45} P.~Deligne, \emph{Cohomologie \'etale (SGA $4\frac12$)}, LNM~569, 1977; \emph{Rapport sur la formule des traces}, 3.2. Background trace-formula reference.
\bibitem{SGA5} A.~Grothendieck (ed.), \emph{Cohomologie \(\ell\)-adique et fonctions L (SGA 5)}, LNM~589, 1977; Expos\'e~III, 4.7. Background Lefschetz--Verdier reference.
\bibitem{SGA6} P.~Berthelot, A.~Grothendieck and L.~Illusie (eds.), \emph{Th\'eorie des intersections et th\'eor\`eme de Riemann--Roch (SGA 6)}, LNM~225, 1971; Expos\'e~V by P.~Berthelot, \S\S2--3. Original scan not visually verified in this run; used as background, not a claimed fresh source-proof check.
\bibitem{FGA} A.~Grothendieck, \emph{Techniques de construction et th\'eor\`emes d'existence en g\'eom\'etrie alg\'ebrique IV: les sch\'emas de Hilbert}, S\'eminaire Bourbaki, Expos\'e~221 (1961), \S3, Theorems~3.1--3.2, pp.~259--260. \url{https://www.numdam.org/article/SB_1960-1961__6__249_0.pdf}.
\bibitem{AK} A.~B.~Altman and S.~L.~Kleiman, \emph{Foundations of the theory of Fano schemes}, Compositio Math.~34 (1977), pp.~3--47; \S2.10, used only to corroborate the SGA~6 locator. \url{https://www.numdam.org/item/CM_1977__34_1_3_0/}.
\bibitem{Archive} Universit\'e de Montpellier, \emph{Archives Grothendieck}, cotes~16, 90, 97, 133, scans specified in \S\ref{sec:archive}. Catalogue: \url{https://grothendieck.umontpellier.fr/archives-grothendieck/}; original files \url{https://grothendieck.umontpellier.fr/16.pdf}, \url{https://grothendieck.umontpellier.fr/90.pdf}, \url{https://grothendieck.umontpellier.fr/97.pdf}, \url{https://grothendieck.umontpellier.fr/133.pdf}. Cote~1 requested leaves remain unverified.
\bibitem{Fable} Claude Fable 5.2, effort max as reported in its supplied file, \emph{FABLE\_PROGRAM\_SUGGESTIONS\_20261120.md} and \emph{\ldots\_EXPLAINED.md}, suggestions~3, 4 and~8. Unchecked leads, not cited theorems.
\bibitem{Richarz} T.~Richarz, \emph{A new approach to the geometric Satake equivalence}, Documenta Mathematica 19 (2014), 209--246. Theorem~3.1, Proposition~4.1, Lemma~5.4, and their proofs; the numbering differs from arXiv 1207.5314v2. \url{https://ems.press/content/serial-article-files/26224?nt=1}.
\bibitem{MV} I.~Mirkovi\'c and K.~Vilonen, \emph{Geometric Langlands duality and representations of algebraic groups over commutative rings}, Annals of Mathematics 166 (2007), 95--143; arXiv math/0401222v5 includes the later Appendix~B correction. Theorem~7.3, Lemma~7.1, Corollary~7.4. \url{https://arxiv.org/abs/math/0401222}.
\bibitem{StacksNil} The Stacks Project, Proposition~59.45.4, tag~03SI, topological invariance of the \'etale topos. \url{https://stacks.math.columbia.edu/tag/03SI}.
\bibitem{Run} This run, \texttt{GCT\_NSRH\_20261121}; findings A01 and F01--F17, with agent source memoranda, proof reviews, computations and \texttt{CHECK\_EVENTS.tsv} in the final package. This is the only current-run identifier used.
\end{thebibliography}
\end{document}
